\documentclass[11pt,a4paper]{article}
\usepackage[utf8]{inputenc}
\usepackage[T1]{fontenc}
\usepackage{amsmath,amssymb,amsthm}
\usepackage{mathtools}
\usepackage{geometry}
\usepackage{fancyhdr}
\usepackage{titlesec}
\usepackage{enumitem}
\usepackage{array}
\usepackage{booktabs}
\usepackage{longtable}

\geometry{margin=1in}

\pagestyle{fancy}
\fancyhf{}
\fancyhead[C]{\small Cayley, Collected Mathematical Papers, Vol.~XII}
\fancyfoot[C]{\thepage}
\renewcommand{\headrulewidth}{0.4pt}

\theoremstyle{definition}
\newtheorem*{art}{Art.}

\DeclareMathOperator{\sech}{sech}
\DeclareMathOperator{\arctanh}{arctanh}

\title{\textbf{The Collected Mathematical Papers of Arthur Cayley}\\[6pt]
\Large Volume XII, Pages 331--380\\[4pt]
\large Papers 845--848}
\author{Arthur Cayley}
\date{}

\begin{document}
\maketitle
\thispagestyle{empty}

\tableofcontents
\newpage

%==============================================================================
\section{Paper 845}
%==============================================================================

\begin{center}
\large\textbf{845.}

\Large\textbf{ON THE ORTHOMORPHOSIS OF THE CIRCLE INTO THE PARABOLA.}

\smallskip
[From the \textit{Quarterly Journal of Pure and Applied Mathematics}, vol.~xxi. (1886), pp.~107--126.]
\end{center}

\bigskip

The function $\log \tan(45^\circ + \frac{1}{2}\arg.)$ is tabulated, Legendre, \textit{Th\'{e}orie des Fonctions Elliptiques}, t.~II. table iv. pp.~256--259, viz. writing as above $\phi$ for the argument, we have hereby the value of $\Phi = \log \tan(\frac{1}{4}\pi + \frac{1}{2}\phi)$, for every value of $\phi$ from $0^\circ$ to $90^\circ$ at intervals of $30'$ and to 12 decimals, and with fifth differences. Observe that $\phi$ is thus given in degrees and minutes as a circular arc, $\Phi$ as a number; it is convenient to have $\phi$ and $\Phi$ each as arcs, or each as numbers, the conversion being of course at once made by means of a table of the lengths of circular arcs, and I have calculated the Table which I give at the end of the present paper.

I had previously calculated for the correspondence of $AO$, $OB$ and $ACB$ with $AO$, $OB$ and $AGB$, the following values, at irregular intervals suited to the construction of a figure.

\bigskip

\textbf{Circle} \hfill \textbf{Parabola}

\medskip

\begin{tabular}{ccc|ccc}
$\phi$ & $P$ & $Q$ & $x$ & $y$ \\
\hline
$0^\circ$ & $0$ & $0$ & $0$ & $0$ \\
$5^\circ$ & $0.0875$ & $0.0437$ & $0.0077$ & $0.0076$ \\
$10^\circ$ & $0.1763$ & $0.0875$ & $0.0306$ & $0.0309$ \\
$15^\circ$ & $0.2680$ & $0.1317$ & $0.0691$ & $0.0706$ \\
$20^\circ$ & $0.3640$ & $0.1763$ & $0.1232$ & $0.1284$ \\
$25^\circ$ & $0.4663$ & $0.2217$ & $0.1935$ & $0.2068$ \\
$30^\circ$ & $0.5774$ & $0.2679$ & $0.2804$ & $0.3094$ \\
$35^\circ$ & $0.7002$ & $0.3153$ & $0.3850$ & $0.4415$ \\
$40^\circ$ & $0.8391$ & $0.3640$ & $0.5085$ & $0.6109$ \\
$45^\circ$ & $1.0000$ & $0.4142$ & $0.6520$ & $0.8284$ \\
$50^\circ$ & $1.1918$ & $0.4663$ & $0.8176$ & $1.1115$ \\
$55^\circ$ & $1.4281$ & $0.5208$ & $1.0086$ & $1.4877$ \\
$60^\circ$ & $1.7321$ & $0.5774$ & $1.2280$ & $2.0000$ \\
$65^\circ$ & $2.1445$ & $0.6371$ & $1.4806$ & $2.7321$ \\
$70^\circ$ & $2.7475$ & $0.7002$ & $1.7732$ & $3.8476$ \\
$75^\circ$ & $3.7321$ & $0.7673$ & $2.1139$ & $5.7274$ \\
$80^\circ$ & $5.6713$ & $0.8391$ & $2.5130$ & $9.5196$ \\
$85^\circ$ & $11.4301$ & $0.9163$ & $2.9835$ & $20.9515$ \\
$90^\circ$ & $\infty$ & $1.0000$ & $3.5452$ & $\infty$
\end{tabular}

\bigskip

The formulae for the orthomorphosis are
\[
\sqrt{x+iy} = p + iq,
\]
where
\[
P = \tan \tfrac{1}{2}\pi p, \quad Q = \tanh \tfrac{1}{2}\pi q,
\]
and
\[
x = p^2 - q^2, \quad y = 2pq.
\]

The relation between the circle and parabola is established by means of the equation
\[
\phi = \log \tan(\tfrac{1}{4}\pi + \tfrac{1}{2}\Phi),
\]
as before, we have $Q = \tan \tfrac{1}{2}\phi$, and the formula is giving the values of $x$, $y$.

It is clear that we have
\[
P = \tan \tfrac{1}{2}\pi p, \quad Q = \tanh \tfrac{1}{2}\pi q,
\]
and thence
\[
p = \frac{2}{\pi}\arctan P, \quad q = \frac{2}{\pi}\arctanh Q.
\]

Hence, to the circle $x^2 + y^2 = c^2$, corresponds in the parabola the curve where $p + iq = \sqrt{x+iy}$, $P = \tan \frac{1}{2}\pi p$, $Q = \tanh \frac{1}{2}\pi q$. This is a complicated transcendental equation, and I do not see any convenient way of tracing the curve.

The set of curves satisfy the differential equation
\[
\frac{P\,dP + Q\,dQ}{PQ(Q\,dP + P\,dQ)} = \text{const.},
\]
that is,
\[
\frac{dP}{P(1+P^2)} + \frac{dQ}{Q(1-Q^2)} = 0,
\]
where $dP$, $dQ$ are given in terms of $dp$, $dq$ by the equations
\[
\frac{dP}{dp} = \tfrac{1}{2}\pi(1+P^2), \quad \frac{dQ}{dq} = \tfrac{1}{2}\pi(1-Q^2).
\]

We have $y = 2pq$, and thence at the highest points, or summits of the several curves, $q\,dp + p\,dq = 0$. Combining these equations, we have $dP : dQ = -P(1+P^2)q : Q(1-Q^2)p$, and thence as an equation to the locus of the summits in question.

If $p$ and $q$ are small, then putting for a moment $\frac{1}{2}\pi = m$ for shortness, we have $P = mp + \frac{1}{3}m^3p^3$, $Q = mq - \frac{1}{3}m^3q^3$, and the equation becomes
\[
p^2(1 + m^2p^2)(1 + m^2q^2) - q^2(1 - m^2q^2)(1 - m^2p^2) = 0,
\]
that is,
\[
p^2 - q^2 + m^2(p^4 + 2p^2q^2 - q^4) = 0.
\]

Writing this in the form we find
\[
x + \tfrac{1}{4}\pi^2(x^2 + 2y^2) = 0,
\]
or say
\[
y^2 = -\frac{2}{\pi^2}x - \tfrac{1}{2}x^2,
\]
as the locus of the summits in the neighbourhood of the focus $O$, viz. the summits lie all of them, as might have been expected, on the left-hand (or negative side) of the focus.

I have constructed for the parabola, to the scale $1 = 1\frac{1}{2}$ inch, as accurately as the data enable, the figure corresponding to the concentric circles and the radii of the circle.

Resuming the equation $\sqrt{x+iy} = p + iq$, that is, $x = p^2 - q^2$ and $y = 2pq$, we have ($p = \text{const.}$), the curves
\[
y^2 = 4p^2(p^2 - x),
\]
and ($q = \text{const.}$),
\[
y^2 = 4q^2(q^2 + x),
\]
which are two systems of confocal parabolas, cutting each other at right angles; the curves of the former set all of them interior to the given parabola, those of the latter set of course cutting it at right angles.

Corresponding hereto in the circle, writing for a moment $\sqrt{X+iY} = p + iq$, we have whence
\[
p = P(1 - qQ), \quad q = Q(1 + pP);
\]
whence eliminating $P$ and $Q$ successively, we find say, for shortness, these are
\[
p^2 + q^2 - 2mp + 1 = 0, \quad \text{and} \quad p^2 + q^2 - 2nq - 1 = 0.
\]

Introducing the polar coordinates $r$, $\theta$, we have $x + iy = r(\cos\theta + i\sin\theta)$, and thence
\[
p, q = \sqrt{r}\cos\tfrac{1}{2}\theta, \quad \sqrt{r}\sin\tfrac{1}{2}\theta;
\]
the equations thus become
\[
r - 1 - 2m\sqrt{r}\cos\tfrac{1}{2}\theta = 0, \quad r + 1 - 2n\sqrt{r}\sin\tfrac{1}{2}\theta = 0,
\]
which belong to two lima\c{c}ons having each of them $B$ for a node and for a focus, and which of course cut each other at right angles; see fig.~4, which is a mere diagram.

In fact, omitting for convenience the accents, but recollecting always that the curves belong to the figure of the circle, the first equation gives
\[
(r - 1)^2 = 4m^2r\cos^2\tfrac{1}{2}\theta = 2m^2r(1 + \cos\theta) = 2m^2(r + x),
\]
that is,
\[
(r^2 + 1 - 2m^2x)^2 = 4(m^2 + 1)^2r^2.
\]

Transforming to the point $B$ as origin, we must write $x + 1$ for $x$, and then the equation thus becomes
\[
[x^2 + y^2 - 2(m^2 + 1)x + 2(m^2 + 1)]^2 = 4(m^2 + 1)^2(x^2 + y^2 + 2x + 1),
\]
that is, or say
\[
(x^2 + y^2)^2 - 4(m^2 + 1)(x^2 + y^2)x + 4(m^2 + 1)(x^2 - m^2y^2) = 0,
\]
showing that the point $B$ is a crunode, the tangents there being $x = \pm my$.

Writing in the equation $y = 0$, we have $x = 0$, the crunode, and
\[
x = 2\sqrt{m^2 + 1}\{\sqrt{m^2 + 1} - m\};
\]
the product of these two values is $= 4(m^2 + 1)$, that is, $= (P + \frac{1}{P})$, viz. one value is greater, the other less, than $2$. We see also that $2\sqrt{m^2 + 1}\{\sqrt{m^2 + 1} - m\}$, the smaller root, is greater than $1$.

The curve corresponding to a parabola $y^2 = 4p^2(p^2 - x)$ is thus a crunodal lima\c{c}on, the crunode at $B$, and the loop lying wholly within the circle. Moreover, the loop includes within itself the centre of the circle.

The other curve is in like manner
\[
(r^2 + 1 + 2n^2x)^2 = 4(n^2 - 1)r^2,
\]
viz. transforming to the origin $B$, and therefore putting $x - 1$ for $x$, and $r^2 = (x-1)^2 + y^2$, the equation is
\[
[x^2 + y^2 + 2(n^2 - 1)x - 2(n^2 - 1)]^2 = 4(n^2 - 1)^2(x^2 + y^2 - 2x + 1),
\]
that is,
\[
[x^2 + y^2 + 2(n^2 - 1)x]^2 = 4n^2(n^2 - 1)(x^2 + y^2),
\]
or say
\[
(x^2 + y^2)^2 + 4(n^2 - 1)(x^2 + y^2)x - 4(n^2 - 1)(x^2 + n^2y^2) = 0,
\]
viz. the point $B$ is an acnode with the imaginary tangents $x = \pm iny$.

Writing in the equation $y = 0$, we have $x = 0$, the acnode, and
\[
x = -\sqrt{n^2 - 1}\{\sqrt{n^2 - 1} + n\},
\]
where $\sqrt{n^2 - 1}$ is real and may be taken to be positive; there is thus one root $-\sqrt{n^2 - 1}\{\sqrt{n^2 - 1} + n\}$ which is negative; the other root, say is positive and less than $1$; the curve is thus an acnodal lima\c{c}on, having $B$ for the acnode and cutting the axis outside the circle on the negative side of $B$, and inside the circle between $B$ and $O$.

Taking $E$ as origin, the equation of the circle is $x^2 + y^2 - 2x = 0$, and we hence find for the intersection of the lima\c{c}on with the circle $n^2x = 2(n^2 - 1)$, that is,
\[
x = 2\left(1 - \frac{1}{n^2}\right), \quad y = \pm\frac{2}{n^2}\sqrt{n^2 - 1}\left(1 - \frac{1}{n^2}\right),
\]
values which belong to a real intersection; for $q = 0$, we have $Q = 0$, $n = \infty$, and therefore $(x,y) = (2,0)$, viz. the intersection is at $B$; for $q = \infty$, we have $Q = 1$, $n = 1$, and therefore $(x,y) = (0,0)$, viz. the intersection is at $A$, results which are obviously right. Observe that for the other lima\c{c}on we have $x = 2(1 + \frac{1}{m^2})$, $y = \pm\frac{2}{m^2}\sqrt{m^2 + 1}(1 + \frac{1}{m^2})$, viz. there is no real intersection with the circle.

\bigskip

\begin{center}
\textbf{The Table above referred to.}

\smallskip
$\Phi = \log \tan(45^\circ + \frac{1}{2}\phi)$.
\end{center}

\medskip

\begin{center}
\small
\begin{tabular}{r|r|r|r|r|r|r|r|r|r|r|r|r|r}
\hline
$\phi$ & $\Phi$ & $\phi$ & $\Phi$ & $\phi$ & $\Phi$ & $\phi$ & $\Phi$ & $\phi$ & $\Phi$ & $\phi$ & $\Phi$ & $\phi$ & $\Phi$ \\
\hline
$0^\circ$ & $0.000000$ & $15^\circ$ & $0.268028$ & $30^\circ$ & $0.549306$ & $45^\circ$ & $0.881374$ & $60^\circ$ & $1.316958$ & $75^\circ$ & $2.027589$ & $90^\circ$ & $\infty$ \\
$1^\circ$ & $0.017455$ & $16^\circ$ & $0.286715$ & $31^\circ$ & $0.570302$ & $46^\circ$ & $0.906320$ & $61^\circ$ & $1.347748$ & $76^\circ$ & $2.092714$ & & \\
$2^\circ$ & $0.034913$ & $17^\circ$ & $0.305497$ & $32^\circ$ & $0.591398$ & $47^\circ$ & $0.931697$ & $62^\circ$ & $1.379619$ & $77^\circ$ & $2.163715$ & & \\
$3^\circ$ & $0.052378$ & $18^\circ$ & $0.324380$ & $33^\circ$ & $0.612600$ & $48^\circ$ & $0.957527$ & $63^\circ$ & $1.412631$ & $78^\circ$ & $2.242090$ & & \\
$4^\circ$ & $0.069853$ & $19^\circ$ & $0.343371$ & $34^\circ$ & $0.633914$ & $49^\circ$ & $0.983832$ & $64^\circ$ & $1.446847$ & $79^\circ$ & $2.330478$ & & \\
$5^\circ$ & $0.087343$ & $20^\circ$ & $0.362476$ & $35^\circ$ & $0.655347$ & $50^\circ$ & $1.010639$ & $65^\circ$ & $1.482336$ & $80^\circ$ & $2.431715$ & & \\
$6^\circ$ & $0.104851$ & $21^\circ$ & $0.381703$ & $36^\circ$ & $0.676905$ & $51^\circ$ & $1.037973$ & $66^\circ$ & $1.519177$ & $81^\circ$ & $2.550723$ & & \\
$7^\circ$ & $0.122381$ & $22^\circ$ & $0.401060$ & $37^\circ$ & $0.698595$ & $52^\circ$ & $1.065864$ & $67^\circ$ & $1.557458$ & $82^\circ$ & $2.694193$ & & \\
$8^\circ$ & $0.139938$ & $23^\circ$ & $0.420554$ & $38^\circ$ & $0.720425$ & $53^\circ$ & $1.094342$ & $68^\circ$ & $1.597277$ & $83^\circ$ & $2.871303$ & & \\
$9^\circ$ & $0.157526$ & $24^\circ$ & $0.440194$ & $39^\circ$ & $0.742403$ & $54^\circ$ & $1.123440$ & $69^\circ$ & $1.638747$ & $84^\circ$ & $3.100349$ & & \\
$10^\circ$ & $0.175148$ & $25^\circ$ & $0.459990$ & $40^\circ$ & $0.764538$ & $55^\circ$ & $1.153195$ & $70^\circ$ & $1.681998$ & $85^\circ$ & $3.405200$ & & \\
$11^\circ$ & $0.192811$ & $26^\circ$ & $0.479950$ & $41^\circ$ & $0.786839$ & $56^\circ$ & $1.183645$ & $71^\circ$ & $1.727179$ & $86^\circ$ & $3.826860$ & & \\
$12^\circ$ & $0.210519$ & $27^\circ$ & $0.500085$ & $42^\circ$ & $0.809317$ & $57^\circ$ & $1.214833$ & $72^\circ$ & $1.774474$ & $87^\circ$ & $4.445411$ & & \\
$13^\circ$ & $0.228279$ & $28^\circ$ & $0.520407$ & $43^\circ$ & $0.831983$ & $58^\circ$ & $1.246804$ & $73^\circ$ & $1.824097$ & $88^\circ$ & $5.434910$ & & \\
$14^\circ$ & $0.246096$ & $29^\circ$ & $0.540925$ & $44^\circ$ & $0.854849$ & $59^\circ$ & $1.279606$ & $74^\circ$ & $1.876305$ & $89^\circ$ & $7.494017$ & & \\
\hline
\end{tabular}
\end{center}

\newpage

%==============================================================================
\section{Paper 846}
%==============================================================================

\begin{center}
\large\textbf{846.}

\Large\textbf{A VERIFICATION IN REGARD TO THE LINEAR TRANSFORMATION OF THE THETA-FUNCTIONS.}

\smallskip
[From the \textit{Quarterly Journal of Pure and Applied Mathematics}, vol.~xxi. (1886), pp.~77--84.]
\end{center}

\bigskip

The notation is that of Smith's ``Memoir on the Theta and Omega Functions,'' differing from Hermite's only in a factor,
\[
\vartheta_{\mu,\nu}(\text{Smith}) = i^{\frac{1}{2}\mu\nu}\vartheta_{\mu,\nu}(\text{Hermite});
\]
and it is in consequence of this that the factor $i^{\frac{1}{2}\mu\nu}$ occurs in the expression presently given for $\delta$. Hermite's Memoir ``Sur quelques formules relatives a la transformation des fonctions elliptiques'' appeared in Liouville, t.~III. (1858), pp.~27--36.

Writing
\[
\omega = \frac{a\Omega + b}{c\Omega + d},
\]
where $ad - bc = 1$, the formula of transformation is
\[
\vartheta_{\mu',\nu'}(\omega) = \delta \cdot \vartheta_{\mu,\nu}(\Omega) \exp\left\{\frac{\pi i}{4}\cdot\frac{a\Omega + b}{c\Omega + d}(c\mu' + d\nu')^2\right\},
\]
where
\[
\mu' = a\mu + b\nu + ab, \quad \nu' = c\mu + d\nu + cd.
\]

We have, according as $b$ is positive or negative,
\[
\delta = \frac{1}{\sqrt{\pm i(c\Omega + d)}} \cdot \frac{H}{\sqrt{(-1)^{\frac{1}{2}(a-1)}}},
\]
where $\left(\frac{a}{b}\right)$, $\left(\frac{c}{d}\right)$, \&c., are the Legendrian symbols as generalised by Jacobi; if $a$ and $b$ are each of them odd, the formulae for the cases $a$ odd and $b$ odd respectively are equivalent to each other, and either may be used indifferently.

To fix the ideas, I assume throughout that $b$ is positive and odd; the proper formulae thus are as above, and
\[
H = \left(\frac{a}{b}\right)i^{-\frac{1}{2}(b-1)}.
\]

I will also, for greater convenience, assume that $c$ is odd; $ad$ is consequently even, viz. the numbers $a$ and $d$ are each of them even, or else one is odd and the other even.

The equation $\omega = \dfrac{a\Omega + b}{c\Omega + d}$ gives $\Omega = \dfrac{-d\omega + b}{c\omega - a}$, and thence
\[
\Omega = \frac{-d + b\omega}{c\omega - a}, \quad H' = \left(\frac{-d}{b}\right)i^{-\frac{1}{2}(b-1)}.
\]

Comparing the expressions for $\omega$, $H$, it appears that we may interchange $\omega$ and $\Omega$, writing also $-d$, $b$, $c$, $-a$ for $a$, $b$, $c$, $d$; and changing the other letters, we may for $a$, $b$, $c$, $d$, $\omega$, $\Omega$, $\mu$, $\nu$, $m$, $n$ write $d$, $b$, $c$, $a$, $\Omega$, $\omega$, $m'$, $n'$, $m$, $n$, where
\[
m' = dm + bn - bd, \quad n' = cm - an - ac.
\]

The formula thus becomes
\[
\vartheta_{m,n}(\Omega) = \delta' \cdot \vartheta_{m',n'}(\omega) \exp\left\{\frac{\pi i}{4}\cdot\frac{d\omega + b}{c\omega - a}(cm' + an')^2\right\},
\]
or, for $x$ writing $(a + b\Omega)x$,
\[
\vartheta_{m',n'}\left(\frac{-d + b\omega}{c\omega - a}\right) = \delta' \cdot \vartheta_{m,n}(\omega) \exp\left\{-\frac{\pi i}{4}\cdot\frac{(c\omega - a)}{(d\omega + b)}(cm' + an')^2\right\},
\]
or, observing that the left-hand side is $= i^{\frac{1}{2}m'n'}\vartheta_{m',n'}$, we have
\[
\vartheta_{m',n'} = \delta' \cdot \vartheta_{m,n} \exp\left\{-\frac{\pi i}{4}\cdot\frac{(c\omega - a)}{(d\omega + b)}(cm' + an')^2\right\},
\]
an equation which is of the same form as the original equation of transformation, and is to be identified with it.

We have
\begin{align*}
m' &= dm + bn - bd = d(a\mu + b\nu + ab) + b(c\mu + d\nu + cd) - bd \\
&= \mu + bd(-a + c - 1),\\
n' &= cm - an - ac = c(a\mu + b\nu + ab) - a(c\mu + d\nu + cd) - ac \\
&= \nu + ac(-b + d - 1),
\end{align*}
values which may be written
\[
m' = 2P - \mu, \quad n' = 2Q - \nu,
\]
where $P = \frac{1}{2}bd(-a + c - 1)$, $Q = \frac{1}{2}ac(-b + d - 1)$, $P$ and $Q$ being integers.

In fact, $P$ is an integer if $b$ or $d$ is even, and if they are each of them odd, then from the equation $ad - bc = 1$, $a$ and $c$ will be one of them odd, the other even, and $a + c - 1$ will be even; and similarly for $Q$.

The left-hand side is
\[
i^{\frac{1}{2}m'n'}\vartheta_{m',n'} = i^{\frac{1}{2}(2P-\mu)(2Q-\nu)}\vartheta_{2P-\mu,\,2Q-\nu},
\]
which is
\[
i^{2PQ - P\nu - Q\mu + \frac{1}{2}\mu\nu}\vartheta_{-\mu,\,-\nu} = i^{2PQ - P\nu - Q\mu} \cdot i^{\frac{1}{2}\mu\nu}\vartheta_{\mu,\nu},
\]
and the equation finally is
\[
\vartheta_{\mu,\nu} = \delta\delta' \cdot \vartheta_{\mu,\nu} \exp\left\{-\frac{\pi i}{4}\cdot\frac{(c\omega - a)}{(d\omega + b)}(cm' + an')^2 + \frac{\pi i}{4}\omega(c\mu + d\nu)^2\right\}.
\]

Comparing this with the original equation the two equations will be identical if only $\delta\delta' = i^{2PQ - P\nu - Q\mu}$.

We have
\[
\delta = \frac{1}{\sqrt{+i(c\Omega + d)}} \cdot \frac{H}{\sqrt{(-1)^{\frac{1}{2}(a-1)}}}, \quad \delta' = \frac{1}{\sqrt{+i(c\omega - a)}} \cdot \frac{H'}{\sqrt{(-1)^{\frac{1}{2}(-d-1)}}},
\]
where the square roots are taken in such wise that the real part is positive; hence
\[
\delta\delta' = \frac{1}{\sqrt{-(c\Omega + d)(c\omega - a)}} \cdot \frac{HH'}{\sqrt{(-1)^{\frac{1}{2}(a-d-2)}}},
\]
where $\sqrt{+1}$ means $\sqrt{-i(a + b\Omega)}$; the last-mentioned two square roots being as just explained; and we have moreover
\begin{align*}
\delta &= \exp\left\{-\tfrac{1}{4}i\pi(ac\mu^2 + 2bc\mu\nu + bd\nu^2 + 2abc\mu + 2abd\nu + ab^2c)\right\}i^{\frac{1}{4}\mu\nu - \frac{1}{2}mn},\\
\delta' &= \exp\left\{-\tfrac{1}{4}i\pi(-dc m^2 + 2bc\,mn - ab n^2 - 2dbcm + 2abdn - db^2c)\right\}i^{\frac{1}{4}m'n'},
\end{align*}
viz. the value of $\delta'$ is obtained from that of $\delta$ by the change $a$, $b$, $c$, $d$, $\mu$, $\nu$, $m$, $n$ into $-d$, $b$, $c$, $a$, $m'$, $n'$, $m$, $n$.

Representing these by
\[
\delta = \exp\{-\tfrac{1}{4}(i\pi)A\}\,i^{\frac{1}{4}\mu\nu - \frac{1}{2}mn}, \quad \delta' = \exp\{-\tfrac{1}{4}(i\pi)A'\}\,i^{\frac{1}{4}m'n'},
\]
we have
\[
\delta\delta' = \exp\{-\tfrac{1}{4}i\pi(A + A')\}\,i^{\frac{1}{4}\mu\nu - \frac{1}{2}mn + \frac{1}{4}m'n'}.
\]

But
\[
m'n' = (2P - \mu)(2Q - \nu) = 4PQ - 2P\nu - 2Q\mu + \mu\nu,
\]
that is,
\[
\tfrac{1}{4}\mu\nu - \tfrac{1}{2}mn + \tfrac{1}{4}m'n' = -PQ + \tfrac{1}{2}P\nu + \tfrac{1}{2}Q\mu,
\]
or, omitting the term divisible by 4,
\[
i^{\frac{1}{4}\mu\nu - \frac{1}{2}mn + \frac{1}{4}m'n'} = i^{\frac{1}{2}P\nu + \frac{1}{2}Q\mu} = \left(\frac{-1}{i}\right)^{\frac{1}{2}P\nu + \frac{1}{2}Q\mu}.
\]

To calculate $A + A'$, we have
\[
dm + bn = \mu + bd(-a + c), \quad cm + an = \nu + ac(-b + d),
\]
and thence
\begin{align*}
&-cdm^2 + (ad + bc)mn - abn^2 =\\
&\quad \mu\nu + \mu\,ac(-b + d) + \nu\,bd(a - c) + abcd(d - b)(a - c) - (ad + bc)\mu\nu,
\end{align*}
also
\begin{align*}
2bcdm &= 2abcd\mu - 2b^2cd\nu - 2ab^2cd,\\
2abdn &= 2abcd\mu + 2abd^2\nu + 2abcd^2,\\
-db^2c &= -db^2c;
\end{align*}
whence, adding, we obtain
\begin{align*}
A' = -ac\mu^2 - 2bc\mu\nu - bd\nu^2 - 2abc\mu + (-2bcd - 2b^2cd + 2abd^2)\nu\\
+ abcd(2bc - ab - cd - 2b + 2d) - db^2c,
\end{align*}
and, adding to this the expression of $A$, we find
\[
A + A' = (2abd - 2bcd - 2b^2cd + 2abd^2)\nu + ab^2c - db^2c + abcd(2bc - ab - cd - 2b + 2d).
\]

The coefficient of $\nu$ is $= 2bd(a - c - bc + ad) = 2bd(a - c + 1)$, which is $= 4P\nu$. Hence writing
\[
\Theta = \tfrac{1}{4}abcd(2bc - ab - cd - 2b + 2d) + \tfrac{1}{2}b^2c(a - d),
\]
where observe that $4\Theta$, but not in every case $\Theta$ itself, is an integer, we have $A + A' = 4P\nu + 4\Theta$, and consequently
\[
\delta\delta' = \left(\frac{-1}{i}\right)^{\Theta},
\]
or, omitting the even number $2P\nu$,
\[
\delta\delta' = \left(\frac{-1}{i}\right)^{\Theta}.
\]

Observe that $\left(\frac{-1}{i}\right)^{\Theta}$ denotes, and it might properly have been written, $\exp\frac{1}{2}i\pi\Theta$. The foregoing equation becomes thus
\[
\delta\delta' = \left(\frac{-1}{i}\right)^{\Theta},
\]
that is,
\[
\delta\delta' = i^{-\Theta},
\]
where the even term $2Q\mu$ may be omitted.

We have moreover
\[
mn = ac\mu^2 + (ad + bc)\mu\nu + bd\nu^2 + ac(b + d)\mu + bd(a + c)\nu + abcd,
\]
\[
\mu\nu = (ad + bc)\mu\nu,
\]
and thence
\begin{align*}
mn - \mu\nu &= ac(\mu^2 - \mu) + 2bc\mu\nu + bd(\nu^2 - \nu) + ac(b + d + 1)\mu\\
&\quad + bd(a + c + 1)\nu + abcd,
\end{align*}
where each term is even; hence $mn - \mu\nu$ is even, and we have simply
\[
\delta\delta' = \left(\frac{-1}{i}\right)^{M},
\]
where
\[
M = \tfrac{1}{4}abcd(ab + cd - 2bc + 2b - 2d) - \tfrac{1}{2}b^2c(a - d).
\]

I write $M = \frac{1}{2}b(a - d)$, then
\[
-M = \tfrac{1}{4}abcd(ab + cd - 2bc + 2b - 2d) - \tfrac{1}{4}(bc + 1)b(a - d),
\]
where the second term is $= \frac{1}{4}abd(a - d)$, and we have therefore
\[
-M = \tfrac{1}{4}abd(abc + c^2d - 2bc^2 + 2bc - 2cd - a + d).
\]

I assume, as above, that $b$ and $c$ are each of them odd; therefore $ad$ is even. I suppose, first, that $ad$ divides by 4, then $\frac{1}{4}abd$ is an integer, and in the expression of $M$, omitting even numbers, we have
\[
abc + c^2d - a + d,
\]
which, putting therein $bc = ad - 1$, becomes
\begin{align*}
&= \tfrac{1}{4}abd\{a^2d + c^2d - 2a + d\},\\
&= \tfrac{1}{4}abd\{a^2d + (c^2 - c)d + (c + 1)d - 2a\},
\end{align*}
where inside the $\{\}$ each term is even; hence $M$ is even.

Next, if $ad$ is even but not divisible by 4, then $bc = ad - 1$, which is $= 1$ (mod.~4), thus $b$ and $c$ are $= 4\sigma + 1$ and $4\tau + 1$, or else $= 4\sigma - 1$ and $4\tau - 1$, and, moreover, $bc = 4\phi + 1$, and $c^2 = 4\phi + 1$; hence
\[
-M = \tfrac{1}{4}abd\{a^2d + (4\phi + 1)d - 2a + d\},
\]
or, omitting even numbers,
\[
-M = \tfrac{1}{4}abd\{a^2d + d - 2a\} = \tfrac{1}{4}abd\{a^2d + (c^2 - c)d + (c + 1)d - 2a\},
\]
that is, inside the $\{\}$ numbers which contain the factor 4, this is
\[
\tfrac{1}{4}abd \cdot 4\{\cdots\} = abd\{\cdots\},
\]
or, since each term within the $\{\}$ is even, we have in this case also $-M$ even.

And this being so, the foregoing equation for $HH'$ becomes
\[
HH' = \left(\frac{a}{b}\right)\left(\frac{-d}{b}\right)i^{-\frac{1}{2}(b-1)}i^{-\frac{1}{2}(b-1)}(-1)^M,
\]
or say
\[
HH' = \left(\frac{a}{b}\right)\left(\frac{-d}{b}\right)(-1)^{\frac{1}{2}(b-1)}.
\]

The values of $H$, $H'$ refer each of them to a positive odd value of $b$, and they thus are
\[
H = \left(\frac{a}{b}\right)i^{-\frac{1}{2}(b-1)}, \quad H' = \left(\frac{-d}{b}\right)i^{-\frac{1}{2}(b-1)};
\]
hence
\[
HH' = \left(\frac{a}{b}\right)\left(\frac{-d}{b}\right)(-1)^{\frac{1}{2}(b-1)},
\]
or, since
\[
\left(\frac{a}{b}\right)\left(\frac{-d}{b}\right) = \left(\frac{-ad}{b}\right) = \left(\frac{1 - bc}{b}\right) = \left(\frac{1}{b}\right) = 1,
\]
and $\frac{1}{2}(b-1)$ divides by 4, this is
\[
HH' = \left(\frac{-ad}{b}\right) = 1.
\]

Also
\[
\frac{c\Omega + d}{c\omega - a} = \frac{c\cdot\frac{-d+b\omega}{c\omega-a} + d}{c\omega - a} = \frac{-cd + bc\omega + cd - ad}{c\omega - a} = \frac{bc - ad}{c\omega - a}\omega = \frac{-1}{c\omega - a}\omega,
\]
but from the equation $ad - bc = 1$, or $\frac{ad - 1}{b} = c$, we have
\[
\frac{c\Omega + d}{c\omega - a} = \frac{a}{b}\cdot\frac{-d + b\omega}{c\omega - a} + \frac{1}{b} = \frac{a\Omega + 1}{b},
\]
whence
\[
HH' = \left(\frac{a}{b}\right)\left(\frac{-d}{b}\right) = \left(\frac{-ad}{b}\right) = 1.
\]

We have $\omega = x + iy$, where $y$ is positive; hence
\[
a + ib\Omega = a + b(x + iy) = a + bx + iby,
\]
if $a = d + bx$; hence
\[
\frac{1}{i(a + b\Omega)} = \frac{1}{i(a + bx) - by} = \frac{-by - i(a + bx)}{b^2y^2 + (a + bx)^2} = \frac{R(\cos\theta + i\sin\theta)}{-R(\cos\theta - i\sin\theta)} = -1,
\]
where $R$ is positive; and $\cos\theta$ is positive since $y$ is positive, and thus $\theta$ lies between $-\frac{\pi}{2}$ and $\frac{\pi}{2}$.

Hence
\[
\sqrt{-i(-d + b\Omega)} = \sqrt{R}(\cos\tfrac{1}{2}\theta + i\sin\tfrac{1}{2}\theta),
\]
being positive, and thus
\[
\sqrt{-(c\Omega + d)(c\omega - a)} = \sqrt{+1} = +1,
\]
that is, $\sqrt{+1} = +1$; and we thus have, as we should do, $HH' = 1$, the equation which was to be verified.

\newpage

%==============================================================================
\section{Paper 847}
%==============================================================================

\begin{center}
\large\textbf{847.}

\Large\textbf{ON THE THEORY OF SEMINVARIANTS.}

\smallskip
[From the \textit{Quarterly Journal of Pure and Applied Mathematics}, vol.~xxi. (1886), pp.~92--107.]
\end{center}

\bigskip

In my ``M\'{e}moire sur les Hyperd\'{e}terminants,'' Crelle, t.~XXX. (1846), pp.~1--37, [13, 14, 16], I gave an investigation for the number of and relations between the quartic invariants of a given binary quantic: the very same formulae apply to the cubic covariants of a given binary quantic, or what is the same thing to the cubic seminvariants of a given weight, and I propose at present (considering the formulae in this point of view) to further develope the theory in regard to the solution of the systems of linear equations obtained in the memoir in question. But I first reproduce the investigation as it stands: I recall that the notation is the ordinary hyperdeterminant notation: for instance,
\[
12^2UV = (a.c - 2b.b + c.a) = 2(ac - b^2);
\]
and that, when the variables do not disappear, each set $(x_1, y_1)$, $(x_2, y_2)$, \ldots\ is ultimately replaced by $(x, y)$, so that these are the variables of any covariant.

The investigation is as follows:

``\ldots We pass to the derivatives of the fourth degree, considering the forms in which all the differential coefficients are of the same order. It is easy to see that we may write
\[
\partial^4 UVWX = (12 \cdot 34)^\alpha (13 \cdot 42)^\beta (14 \cdot 23)^\gamma \cdot UVWX,
\]
or $D_{\alpha,\beta,\gamma}$.''

\medskip

Writing for shortness
\[
12 \cdot 34 = 21, \quad 13 \cdot 42 = 33, \quad 14 \cdot 23 = \mathfrak{S},
\]
we have
\[
21 + 33 + \mathfrak{S} = 0.
\]

Suppose $U = V = W = X$, and consider the derivatives which correspond to the same value $f$ of $\alpha$, $\beta$, $\gamma$: we have to find how many of these derivatives are independent, and to express the others in terms of these. Since the functions are equal after the differentiations, we may before the differentiations interchange in any manner the symbolic numbers $1$, $2$, $3$, $4$: this gives
\[
D_{\alpha,\beta,\gamma} = D_{\alpha,\gamma,\beta} = D_{\beta,\alpha,\gamma} = D_{\beta,\gamma,\alpha} = D_{\gamma,\alpha,\beta} = D_{\gamma,\beta,\alpha}.
\]

But we have identically
\[
(21 + 33 + \mathfrak{S}) \cdot 21^{\alpha-1} 33^\beta \mathfrak{S}^\gamma \cdot UVWX = 0,
\]
when $\alpha > 0$. Multiplying by $21^a 33^b \mathfrak{S}^c$, and applying it to the product $UVWX$, we have
\[
D_{a+1,\,b,\,c} + D_{a,\,b+1,\,c} + D_{a,\,b,\,c+1} = 0,
\]
which, putting $a + b + c = f$, gives a system of equations between the derivatives $D_{\alpha,\beta,\gamma}$ for which $\alpha + \beta + \gamma = f$.

Reducing these by the conditions just obtained, suppose that $f$ denotes the number of partitions of $f$ into three parts, zero included, and the permutations of the parts being disregarded, we have a number $f$ of derivatives $D_{\alpha,\beta,\gamma}$, and a number $\theta(f - 1)$ of linear relations between these derivatives. There remains therefore a number $f - \theta(f - 1)$ of independent derivatives: but when $f$ is an even number, we have included among these the functions $D_{\frac{1}{2}f,\,0,\,0}$, that is, $12^{\frac{1}{2}f} \cdot 34^{\frac{1}{2}f}UVWX$, which evidently reduces itself to $12^{\frac{1}{2}f}U \cdot 34^{\frac{1}{2}f}X$, or $B_{\frac{1}{2}f}(U,V) \cdot B_{\frac{1}{2}f}(W,X)$, that is to say, to the square of $B_{\frac{1}{2}f}(U,V)$.

We must therefore diminish by unity this number $f - \theta(f - 1)$, when $f$ is even. Let $E(\frac{1}{2}f)$ denote the integer part of the fraction $\frac{1}{2}f$: it can be shown that the required number is equal to $E(\frac{1}{2}f) - E(\frac{1}{3}f)$ according as $f$ is even or odd.

Giving to $f$ the six forms $6g$, $6g+1$, $6g+2$, $6g+3$, $6g+4$, $6g+5$, we obtain for the independent derivatives of the fourth degree the corresponding numbers
\[
0,\quad 0,\quad 0,\quad g+1,\quad 0,\quad g+1;
\]
there is for example a single derivative for each of the orders $3$, $5$, $6$, $7$, $8$, $10$, two for each of the orders $9$, $11$, $12$, $13$, $14$, $16$, \&c.

We may take for independent derivatives, when $f$ is even, the terms of the series
\[
D_{f-3,\,3,\,0},\quad D_{f-6,\,6,\,0},\quad \ldots,
\]
and when $f$ is odd, those of the series
\[
D_{f-1,\,1,\,0},\quad D_{f-4,\,4,\,0},\quad D_{f-7,\,7},\quad \ldots,
\]
continuing each series until the last term in which the first suffix exceeds or is equal to the second suffix.

We have in each case the right number of terms; for example, in the case $f = 9$, we take for independent derivatives the two functions $D_{8,1,0}$ and $D_{5,4,0}$, and we form the equations
\begin{align*}
D_{9,0,0} + D_{8,1,0} + D_{7,2,0} &= 0,\\
D_{6,2,1} + D_{5,3,1} + D_{4,4,1} &= 0,\\
D_{7,1,1} + D_{6,2,1} + D_{5,3,1} &= 0,\\
D_{4,3,2} + D_{3,4,2} + D_{2,5,2} &= 0,\\
D_{6,1,2} + D_{5,2,2} + D_{4,3,2} &= 0,\\
D_{3,3,3} + D_{2,4,3} + D_{1,5,3} &= 0,\\
D_{5,1,3} + D_{4,2,3} + D_{3,3,3} &= 0,\\
D_{2,3,4} + D_{1,4,4} + D_{0,5,4} &= 0,\\
D_{4,1,4} + D_{3,2,4} + D_{2,3,4} &= 0,\\
D_{1,3,5} + D_{0,4,5} + D_{?,?} &= 0,
\end{align*}
which are to be reduced by means of the formulas
\[
D_{\alpha,\beta,\gamma} = D_{\alpha,\gamma,\beta} = D_{\beta,\gamma,\alpha} = \cdots
\]

We will presently give the solution of these equations: but beginning with the second order and going successively to the ninth order, we form easily the following table:

\begin{align*}
D_{2,0,0} &= B_2, & D_{3,0,0} &= 0, & D_{4,0,0} &= 0,\\
D_{2,1,0} &= 0, & D_{3,1,0} &= 0, & D_{4,1,0} &= 0,\\
D_{1,1,1} &= 0, & D_{2,2,0} &= B_2^2, & D_{3,2,0} &= 0,\\
& & D_{2,1,1} &= 0, & D_{3,1,1} &= 0,\\
& & & & D_{2,2,1} &= 0,\\
& & & & D_{1,1,2} &= 0.
\end{align*}

\begin{align*}
D_{5,0,0} &= 0, & D_{6,0,0} &= B_2^3, & D_{7,0,0} &= 0,\\
D_{5,1,0} &= 0, & D_{6,1,0} &= 0, & D_{7,1,0} &= 0,\\
D_{4,2,0} &= D_{4,1,1} = 0, & D_{5,2,0} &= \tfrac{1}{3}B_2 D_{4,2,0} = 0, & D_{6,2,0} &= \tfrac{1}{2}B_2 D_{5,2,0} = 0,\\
D_{3,3,0} &= B_2^3, & D_{4,3,0} &= \tfrac{1}{2}B_2^2 D_{3,3,0} = \tfrac{1}{2}B_2^3, & D_{5,3,0} &= \tfrac{1}{5}B_2 D_{4,3,0} = \tfrac{1}{10}B_2^4,\\
D_{3,2,1} &= 0, & D_{4,2,1} &= 0, & D_{5,2,1} &= 0,\\
& & D_{3,3,1} &= B_2^3, & D_{4,3,1} &= \tfrac{1}{4}B_2^4,\\
& & & & D_{3,2,2} &= 0.
\end{align*}

For any value of $f$ except $f = 2$, $3$, or $4$, the table commences, for $f$ even and $f$ odd respectively, in the following manner:
\begin{align*}
D_{f,0,0} &= B_2^{\frac{1}{2}f}, & D_{f,0,0} &= 0,\\
D_{f-1,\,1,\,0} &= -\tfrac{1}{2}B_2^{\frac{1}{2}f}, & D_{f-1,\,1,\,0} &= B_2^{\frac{1}{2}(f-1)},\\
D_{f-2,\,2,\,0} &= -D_{f-1,\,1,\,0}, & D_{f-2,\,1,\,1} &= 0,\\
D_{f-3,\,3,\,0} &= \tfrac{1}{3}B_2 D_{f-2,\,2,\,0} + \tfrac{2}{3}D_{f-1,\,1,\,0},
\end{align*}
but beyond this I do not know the law of the series.''

Before going further, I remark that if $\mathfrak{A}$, $\mathfrak{B}$, $\mathfrak{C}$, instead of the foregoing values, denote respectively
\begin{align*}
\mathfrak{A} &= (00x_1 + 21y_1) \cdot 23 = 0,\\
\mathfrak{B} &= (00x_2 + 21y_2) \cdot 31 = 0,\\
\mathfrak{C} &= (00x_3 + 21y_3) \cdot 12 = 0,
\end{align*}
then identically $\mathfrak{A} + \mathfrak{B} + \mathfrak{C} = 0$, and the same theory applies to the cubic derivatives $\mathfrak{A}^p \mathfrak{B}^q \mathfrak{C}^r \cdot UVW$, that is, to the cubic covariants, or, attending only to the coefficients of the highest powers of $x$, to the cubic seminvariants.

Or, we may use the Clebschian notation, for instance
\[
(a, b, c \mid x, y)^2 = (ax + by)^2 = (a_0x + a_1y)^2 = \text{\&c.},
\]
that is
\[
a_0^2 = a_1^2 = \cdots = a, \quad a_0a_1 = b_0b_1 = \cdots = b, \quad a_1^2 = b_1^2 = \cdots = c,
\]
viz. in the language of Prof. Sylvester, $a_0$, $a_1$, $b_0$, $b_1$, \ldots\ are here umbrae.

The invariant is
\[
(a_0b_1 - a_1b_0)^2 = a_0^2b_1^2 - 2a_0a_1b_0b_1 + a_1^2b_0^2 = ac - 2b \cdot b + ac = 2(ac - b^2),
\]
and so in other cases.

To apply this directly to the theory of seminvariants, it is more convenient to write
\[
(1, b, c \mid x, y)^2 = (x + \alpha y)^2 = \text{\&c.},
\]
that is
\[
\alpha = \beta = \cdots = b,
\]
where $\alpha$, $\beta$, \ldots\ are umbrae.

The invariant is
\[
(\alpha - \beta)^2 = \alpha^2 - 2\alpha\beta + \beta^2 = c - 2b \cdot b + c = 2(c - b^2),
\]
or, to take the case of a cubic function
\[
(1, b, c, d \mid x, y)^3 = (x + \alpha y)^3 = (x + \beta y)^3 = (x + \gamma y)^3,
\]
that is
\[
\alpha = \beta = \gamma = b, \quad \alpha^2 = \beta^2 = \gamma^2 = c, \quad \alpha^3 = \beta^3 = \gamma^3 = d,
\]
a seminvariant is
\begin{align*}
(\alpha - \beta)^2(\alpha - \gamma) &= \alpha^3 - 2\alpha^2\beta + \alpha\beta^2 - \alpha^2\gamma + 2\alpha\beta\gamma - \beta^2\gamma\\
&= d - 2bc + bc - bc + 2b^3 - bc\\
&= d - 3bc + 2b^3,
\end{align*}
and so in other cases.

Writing here the general form of a cubic seminvariant of the weight $\omega$ is now $21^p 33^q \mathfrak{S}^r$, or say $D_{p,q,r}$ where $p + q + r = \omega$ (the new letters $p$, $q$, $r$ being used for the exponents, since $\alpha$, $\beta$, $\gamma$ are now employed in a different sense; and since $f$ will be required as a coefficient, I have also written $\omega$ instead of $f$ for the weight).

We have, as before,
\[
D_{p,q,r} = D_{p,r,q} = D_{q,p,r} = \cdots
\]
and we again see that the theory as to the number of and relations between the cubic seminvariants is identical with that of the quartic invariants.

Observe also that for $\omega$ odd, $D_{\omega,0,0}$ is $= 0$, while for $\omega$ even, $D_{\omega,0,0} = 21^{\omega}$, is a quadric seminvariant, which is of course not to be reckoned among the proper cubic seminvariants; this exactly corresponds to the $B_f$ which is not a proper quartic invariant.

The choice of the functions $D_{f,\,0,\,0}$, $D_{f-3,\,3,\,0}$, \ldots, or $D_{f-1,\,1,\,0}$, $D_{f-4,\,4,\,0}$, \ldots\ for expressing in terms of them the other functions has its advantages, but also its disadvantages; in what follows, I in effect replace them by $D_{f,\,0,\,0}$, $D_{f-2,\,2,\,0}$, \ldots\ and $D_{f-1,\,1,\,0}$, $D_{f-3,\,3,\,0}$, \ldots\ respectively.

And instead of considering the whole series of the functions $D_{p,\,q,\,r}$ for a given weight $\omega$, I consider only those of the form $D_{p,\,q,\,0}$; these form a single series $D_{\omega,0,0}$, $D_{\omega-1,\,1,\,0}$, $D_{\omega-2,\,2,\,0}$, \ldots; and for shortness, I call them $A$, $B$, $C$, $D$, \&c.\ respectively, viz.\ I write
\[
A = D_{\omega,0,0}, \quad B = D_{\omega-1,\,1,\,0}, \quad C = D_{\omega-2,\,2,\,0}, \quad D = D_{\omega-3,\,3,\,0}, \quad \ldots
\]

Suppose first that $\omega$ is odd; we have
\[
D_{p,q,r} = D_{p,r,q} = D_{q,r,p} = \cdots
\]
and, in particular,
\[
D_{\omega-2q,\,q,\,q} = D_{\omega-2q,\,q,\,q},
\]
that is,
\[
D_{\omega-2q,\,q,\,q} = 0,
\]
that is,
\[
C = 0, \quad F = 0, \quad I = 0, \quad \ldots
\]
Hence
\[
B + C = 0, \quad C + 2D + E = 0, \quad D + 3E + 3F + G = 0, \quad \ldots
\]
that is, we have
\[
A = 0, \quad B + C = 0, \quad C + 2D + E = 0, \quad D + 3E + 3F + G = 0, \quad \text{\&c.}
\]
The readiest way of solving these is to express the other functions in terms of $B$, $D$, $F$, $H$, \&c.; viz.\ we thus have
\begin{align*}
A &= 0,\\
C &= -B,\\
E &= B - 2D,\\
G &= 3B - 5D + 3F,\\
I &= -18B + 28D - 14F + 4H,\\
K &= 57B - 81D + 45F - 15H + 5J, \quad \text{\&c.}
\end{align*}

We have here, writing for shortness $900$, $810$, \&c., instead of $D_{9,0,0}$, $D_{8,1,0}$, \&c.\ respectively,
\begin{align*}
900 &= 0, & 810 &= B, & 720 &= -B, & 630 &= D,\\
540 &= B - 2D, & 450 &= F, & 360 &= -3B + 5D - 3F,\\
270 &= H, & 180 &= 18B - 28D + 14F - 4H, & 090 &= L.
\end{align*}
But we have $900 = 0$, $810 + 180 = 0$, $720 + 270 = 0$, $630 + 360 = 0$, $540 + 450 = 0$; that is,
\[
L = 0, \quad 18B - 28D + 14F - 4H = 0, \quad -B + H = 0, \quad -3B + 6D - 3F = 0, \quad B - 2D + F = 0,
\]
all satisfied by $F = B + 2D$, $H = B$, $L = 0$.

The proper course is to stop at the equation for $540$, viz.\ the system is
\[
900 = 0, \quad 810 = B, \quad 720 = -B, \quad 630 = D, \quad 540 = B - 2D,
\]
equations which may be considered as expressing the several functions $900$, $810$, $720$, $630$, $540$ in terms of $B = 810$ and $D = 630$. They agree, as they should do, with a foregoing result, $900 = 0$, $810$, $720 = -810$, $630 = \frac{1}{3}810 + \frac{2}{3}540$, $540$, and it would of course be possible to express the remaining forms $711$, $621$, $522$, $441$, $432$, and $333$ in terms of $B = 810$ and $D = 630$. But observe that the speciality of the present process is that, instead of the whole series of forms $900$, $810$, $720$, $711$, $630$, $621$, $540$, $531$, $522$, $441$, $432$, $330$, we work only with the forms $900$, $810$, $720$, $630$, $540$, for which the last element is $= 0$.

A more simple solution of the system of linear equations in $A$, $B$, $C$, \ldots\ is the following:

\medskip
\begin{center}
\textbf{A B C D E F G H I J K L}
\end{center}
\medskip

For the case $\omega$ even, I write
\[
A + 2B = A', \quad B + 2C = B', \quad C + 2D = C', \quad \text{\&c.;}
\]
and I say that the new functions $A'$, $B'$, $C'$, \ldots\ are related together precisely in the same manner as are the functions $A$, $B$, $C$, \ldots\ belonging to an odd weight $\omega$; viz.\ we have
\[
A' = 0, \quad B' + C' = 0, \quad C' + 2D' + E' = 0,
\]
which being so, the theory is included in what precedes, and there is no occasion to consider it separately.

To prove this, observe that $\omega$ being even, we have
\[
D_{p,q,r} = D_{p,r,q} = D_{q,r,p} = \cdots
\]
that is,
\[
2D_{\omega-2q,\,q,\,q} = D_{\omega-2q+1,\,q,\,q-1} + D_{\omega-2q-1,\,q+1,\,q},
\]
whence, in particular,
\[
2D_{\omega,0,0} = D_{\omega-1,\,1,\,0}, \quad 2D_{\omega-2,\,1,\,1} = D_{\omega-1,\,1,\,0} + D_{\omega-3,\,2,\,1}, \quad \text{\&c.;}
\]
we thus have
\[
A' = 2D_{\omega,0,0} + 2D_{\omega-2,\,2,\,0} = 2D_{\omega-1,\,1,\,0} + D_{\omega-2,\,2,\,0} = 2D_{\omega-1,\,1,\,0}(21 + 33 + \mathfrak{S}),
\]
which is $= 0$.

Similarly $B' + C' = B' + 3C' + \mathfrak{S}' = D_{\omega-3,\,3,\,0} + 3D_{\omega-4,\,4,\,0} + 2D_{\omega-5,\,5,\,0}$, $= D_{\omega-3,\,3,\,0}(2(21 + 33)(21 + 2\cdot 33))$, $= D_{\omega-3,\,3,\,0}(2(\mathfrak{S} + 33)(21 + 2\cdot 33))$, or, since $21 + 33 + \mathfrak{S} = 0$, this is $= -2D_{\omega-3,\,3,\,0}(21 + 33)(21 + 2\cdot 33)$, which, in virtue of $21 + 33 + \mathfrak{S} = 0$, is $= -2D_{\omega-3,\,3,\,0}(21 + 33)(\mathfrak{S} - 21)$, which is $= 0$.

And similarly $C' + 2D' + E' = 0$, \&c.

I come to a new part of the theory: to fix the ideas, consider the weight $\omega = 27$; and write also
\[
a_0, a_1, a_2, a_3, \ldots = 1, b, c, d, e, f, g, h, i, j, k, l, m, \ldots,
\]
using, if we please, the suffixed $a$ for the higher terms.

The cubic seminvariants are
\[
3^{12}, \quad 3^{10}5, \quad 3^85^2, \quad 3^65^3, \quad 3^45^4, \quad \text{viz.\ the number of them is} = 5;
\]
we have thus the 5 terms $B$, $D$, $F$, $H$, $J$.

Here $B = (\alpha - \beta)^{26}(\gamma - \alpha) = (\alpha - \beta)^{26}(\beta - \gamma)$, or, interchanging the $\alpha$ and $\beta$, $= (\alpha - \beta)^{26}(\alpha - \gamma)$, there is an initial term $\alpha^{27} = a^{27}$ and a final term $\alpha\beta^{13}\gamma^{13} = bm^2$, or we may write
\[
B = (\alpha - \beta)^{26}(\alpha - \gamma) = a^{27} + bm^2;
\]
and similarly for the other forms.

We have thus
\begin{align*}
B &= (\alpha - \beta)^{26}(\alpha - \gamma) = a^{27} + bm^2,\\
D &= (\alpha - \beta)^{24}(\alpha - \gamma)^3 = a^{27} + dm^2,\\
F &= (\alpha - \beta)^{20}(\alpha - \gamma)^5 = a^{25}e + fm^2,\\
H &= (\alpha - \beta)^{16}(\alpha - \gamma)^7 = a^{21}g + hm^2,\\
J &= (\alpha - \beta)^{12}(\alpha - \gamma)^9 = a^{19}i + jm^2,
\end{align*}
where the initial term $a^{27} = a^{27}$ has in each case the coefficient unity, but the final terms have each of them the proper numerical coefficient.

It must be possible to form linear combinations
\begin{align*}
B &= a^{27} + bm^2,\\
(B, D) &= a^{25}(c + b^2) + dm^2,\\
(B, D, F) &= a^{23}(e + c^2) + fm^2,\\
(B, D, F, H) &= a^{21}(g + d^2) + hm^2,\\
(B, D, F, H, J) &= a^{19}(i + e^2) + jm^2,
\end{align*}
where $c + b^2$, $e + c^2$, $g + d^2$, $i + e^2$ denote the seminvariants
\[
c - b^2, \quad e - 4bd + 3c^2, \quad g - 6bf + 15ce - 10d^2, \quad i - 8bh + 28cg - 56df + 35e^2,
\]
respectively.

The expressions indicated by their initial and final terms $a^{27} + bn^2$, $a^{25}c + dm^2$, \&c., are what I call ``columns;'' see my paper ``Seminvariant Tables,'' American Journal of Mathematics, t.~VII. (1885), pp.~59--73, [831], where, however, the theory is not by any means completely developed.

As to this, observe that $B$, $D$ are each of them of the form $a^{27} + a^{26}b + a^{25}(c, b^2) + \ldots$, the proper combination in order to eliminate $a^{26}b$ is obviously $B - D$, the term in $a^{26}b$ will then disappear of itself, and the terms in $a^{25}(c, b^2)$ will combine into a term $a^{25}(c - b^2)$, viz.\ in the function qua seminvariant the coefficient of the highest letter $a^{25}$ will be the seminvariant $c - b^2$.

Similarly, in the combination $B$, $D$, $F$, the two coefficients will be determinable so that there are no terms in $a^{27}$, $a^{26}b$, $a^{25}$ or $a^{24}$, and this being so, the term in $a^{23}$ will be $a^{23}(e - 4bd + 3c^2)$, viz.\ in the function, qua seminvariant, the coefficient of the highest letter $a^{23}$ will be the seminvariant $e - 4bd + 3c^2$. And similarly for the remaining two linear combinations.

As a partial verification, I write
\[
B = \{\alpha^{26} + \beta^{26} - 26(\alpha^{25}\beta + \alpha\beta^{25}) + 325(\alpha^{24}\beta^2 + \alpha^2\beta^{24}) \ldots\}(\alpha - \gamma),
\]
or, retaining only the terms which have an index at least $= 25$, this is
\begin{align*}
B &= \alpha^{27} + \alpha\beta^{26} - 26(\alpha^{26}\beta + \alpha^2\beta^{25}) + 325\alpha^{25}\beta^2\\
&\quad - (\alpha^{26}\gamma - 26\alpha^{25}\beta\gamma + 325\alpha^{24}\beta^2\gamma)\\
&= \alpha^{27} - 27\alpha^{26}\beta + \alpha^{25}(299c + 325b^2);
\end{align*}
and similarly
\begin{align*}
D &= (\alpha - \beta)^{24}(\alpha - \gamma)^3\\
&= (\alpha^{24} + \beta^{24} - 24(\alpha^{23}\beta + \alpha\beta^{23}) + 276(\alpha^{22}\beta^2 + \alpha^2\beta^{22}) \ldots\}\\
&\quad \times (\alpha^3 - 3\alpha^2\gamma + 3\alpha\gamma^2 - \gamma^3),
\end{align*}
or, retaining only the terms which have an index at least $= 25$,
\begin{align*}
D &= \alpha^{27} - 24\alpha^{26}\beta + 276\alpha^{25}\beta^2 - 3\alpha^{26}\gamma + 72\alpha^{25}\beta\gamma\\
&\quad + 3\alpha^{24}\beta^2\gamma + 276\alpha^{25}\beta^2 + 72\alpha^{25}\beta\gamma\\
&\quad + 325(279c + 72b^2),
\end{align*}
and thence
\[
B - D = \alpha^{25}(20c - 20b^2),
\]
viz.\ the coefficient of $a^{25}$ is $= 20(c - b^2)$.

But instead of the foregoing forms $B$, $D$, $F$, $H$, $J$, we may start with five other forms which are in fact linear combinations of these, viz.\ these are the forms
\begin{align*}
X &= (\alpha - \beta)^{26}(\alpha + \beta - 2\gamma) = a^{27} + bm^2,\\
Y &= (\alpha - \beta)^{24}(\alpha + \beta - 2\gamma)(\alpha - \gamma)(\beta - \gamma) = a^{25}(c + b^2) + dm^2,\\
Z &= (\alpha - \beta)^{20}(\alpha + \beta - 2\gamma)(\alpha - \gamma)^2(\beta - \gamma)^2 = a^{23}(e + c^2) + fm^2,\\
W &= (\alpha - \beta)^{16}(\alpha + \beta - 2\gamma)(\alpha - \gamma)^3(\beta - \gamma)^3 = a^{21}(g + d^2) + hm^2,\\
T &= (\alpha - \beta)^{12}(\alpha + \beta - 2\gamma)(\alpha - \gamma)^4(\beta - \gamma)^4 = a^{19}(i + e^2) + jm^2,
\end{align*}
where as well the initial as the final terms should have their proper numerical coefficients.

As to this, observe that in $Y$ we have a term $\alpha^{26}(\beta - \gamma)$ which is $= 0$, and similarly a term $\beta^{26}(\alpha - \gamma)$ which is $= 0$; the highest powers are thus $\alpha^{25}$ and $\beta^{25}$, giving a term in $a^{25}$ which can only be $a^{25}(c + b^2)$.

In $Z$ we have the terms $\alpha^{25}(\beta - \gamma)^2$ and $\beta^{25}(\alpha - \gamma)^2$, giving the term $a^{25}(c + b^2)$; in $W$ we have the terms $\alpha^{24}(\beta - \gamma)^3$ and $\beta^{24}(\alpha - \gamma)^3$ which are each $= 0$, the highest powers thus are $\alpha^{23}$, $\beta^{23}$ giving a term in $a^{23}$ which can only be $a^{23}(e - 4bd + 3c^2)$; and in $T$ we have the terms $\alpha^{23}(\beta - \gamma)^4$ and $\beta^{23}(\alpha - \gamma)^4$, which give the term $a^{23}(e - 4bd + 3c^2)$.

The combinations which give the columns thus are
\begin{align*}
X &= a^{27} + bm^2,\\
(Y) &= a^{25}(c + b^2) + dm^2,\\
(Y, Z) &= a^{23}(e + c^2) + fm^2,\\
(Y, Z, W) &= a^{21}(g + d^2) + hm^2,\\
(Y, Z, W, T) &= a^{19}(i + e^2) + jm^2,
\end{align*}

I start with the form $u = (Y, Z, W, T) = a^{19}(i + e^2) + jm^2$.

We know that it is possible to determine the coefficients $\lambda$, $\mu$, $\nu$, in suchwise that the linear function $u + \lambda Z + \mu W + \nu T$ shall be of the proper form $a^{19}(i + e^2) + jm^2$; or, what is the same thing, that there shall be no term with an index exceeding $19$; the conditions to be satisfied are apparently more than three, but of course the number of independent conditions must be $= 3$.

We have, in particular, to get rid of all the terms of the second degree which precede $a^{19}i$, viz.\ the coefficients of $a^{26}c$, $a^{24}d$, $a^{22}e$, $a^{20}f$, $a^{18}g$, $a^{16}h$ must all of them vanish.

Now the terms of $u$ which produce terms of the second degree are the terms containing any two but not all three of the symbolic letters $\alpha$, $\beta$, $\gamma$, say the biliteral terms; we have for instance $\alpha^{25}\beta^2$, $\alpha^{25}\gamma^2$, $\beta^{25}\alpha^2$, $\beta^{25}\gamma^2$ (there are no terms $\gamma^{25}\alpha^2$ or $\gamma^{25}\beta^2$, since the highest power of $\gamma$ is $\gamma^9$), each of which is $= a^{25}c$. But in the development of $u$, we may in any biliteral term by an interchange of the letters $\alpha$, $\beta$, $\gamma$ make the letter of highest index to be $\alpha$, and the other letter to be $\beta$; thus the several terms $\alpha^{25}\beta^2$, $\alpha^{25}\gamma^2$, $\beta^{25}\alpha^2$, $\beta^{25}\gamma^2$ may be each of them converted into $\alpha^{25}\beta^2$, and so in other cases.

Imagining this to be done, the term in $a^{25}c$ is simply the term in $\alpha^{25}\beta^2$, and in like manner the term in $a^{24}d$ is the term in $\alpha^{24}\beta^3$, and so in other cases; the condition as to the disappearance of the terms $a^{25}c$, \ldots, $a^{16}h$ is the condition that the terms in $\alpha^{25}\beta^2$, $\alpha^{24}\beta^3$, $\alpha^{23}\beta^4$, $\alpha^{22}\beta^5$, $\alpha^{21}\beta^6$, $\alpha^{20}\beta^7$ shall all of them disappear.

And if, in the function $u$ transformed in the foregoing manner, we write for convenience $\alpha = 1$, so that $u$ has become a function of $\beta$ only (and observe that $u$ will contain the factor $\beta^2$), then the condition is that the terms in $\beta^2$, $\beta^3$, $\beta^4$, $\beta^5$, $\beta^6$, $\beta^7$ shall all of them disappear.

The conditions may in fact be written
\[
u = 0; \quad \text{viz.\ it is assumed that $u$ is in the first instance transformed into a function of $\beta$ as just explained,}
\]
and the equation is then to be understood as denoting that the coefficients of $u$ are to be so determined that as many as possible of the terms (beginning with that which contains the lowest power of $\beta$) shall each of them vanish.

Consider any term of $u$, for instance the term
\[
T = (\alpha - \beta)^{12}(\alpha + \beta - 2\gamma)(\alpha - \gamma)^4(\beta - \gamma)^4.
\]

To obtain the biliteral terms in the proper form
\begin{enumerate}[label=\arabic*.]
\item we write therein $\gamma = 0$,
\item we write $\beta = 0$, changing $\gamma$ into $\beta$;
\item we write $\alpha = 0$, changing $\beta$, $\gamma$ into $\alpha$, $\beta$;
\end{enumerate}
we thus obtain
\begin{align*}
&(\alpha - \beta)^{12} \alpha^{12} (-2\beta)(-\alpha)^4(-\beta)^4,\\
&\quad (\alpha - \beta)^{12} \alpha^{12} (-2\beta)(-\alpha)^4(-\beta)^4,
\end{align*}
viz.\ the second and third terms are identical; the first term requires, however, the factor $2$ (for any term $\alpha^p\beta^q$ is accompanied by a corresponding term $\alpha^q\beta^p$), so that, omitting this factor throughout, the terms are
\[
(\alpha - \beta)^{24}(\alpha + \beta)\alpha\beta - \alpha^{24}(\alpha - 2\beta)(\alpha - \beta)\beta,
\]
or, putting therein $\alpha = 1$, the terms are
\[
(1 - \beta)^{24}(1 + \beta)\beta - (1 - 2\beta)(1 - \beta)\beta.
\]

Similarly from
\[
Z = (\alpha - \beta)^{20}(\alpha + \beta - 2\gamma)(\alpha - \gamma)^2(\beta - \gamma)^2,
\]
we have
\[
(1 + \beta)\beta(1 - \beta)^{20} + \lambda(1 - \beta)^{20}\beta^2 + \mu(1 - \beta)^{16}\beta^3 + \nu(1 - \beta)^{12}\beta^4
\]
and the like as regards $W$ and $T$.

The result thus is
\begin{align*}
&(1 + \beta)\beta(1 - \beta)^{24} + \lambda(1 - \beta)^{22}\beta^2 + \mu(1 - \beta)^{20}\beta^3 + \nu(1 - \beta)^{18}\beta^4\\
&\quad - (1 - 2\beta)(1 - \beta)\beta\{1 - \lambda\beta(1 - \beta) + \mu\beta^2(1 - \beta)^2 - \nu\beta^3(1 - \beta)^3\} = 0,
\end{align*}
or, as this may be written,
\[
(1 - \beta)^{24}(1 + \beta)\beta = (1 - 2\beta)(1 - \beta)\beta\{1 - \lambda\beta(1 - \beta) + \mu\beta^2(1 - \beta)^2 - \nu\beta^3(1 - \beta)^3\},
\]
viz.\ each side is to be expanded in ascending powers of $\beta$, and the coefficients $\lambda$, $\mu$, $\nu$ are then to be determined so that as many terms as possible shall vanish.

Expanding, we have
\begin{align*}
&(1 - 20\beta + 190\beta^2 - 1138\beta^3 + 4808\beta^4 - 15178\beta^5 + \ldots)\\
&\quad \times \{1 + \lambda\beta + \beta^2(2\lambda + \mu) + \beta^3(3\lambda + \nu) + \beta^4(4\lambda + 10\mu + 6\nu)\\
&\qquad\quad + \beta^5(5\lambda + 20\mu + 21\nu) + \ldots\}\\
&\quad - 1 + \lambda\beta + \beta^2(-\lambda - \mu) + \beta^3(2\mu + \nu) + \beta^4(-3\nu) + \beta^5 \cdot 3\nu + \ldots = 0,
\end{align*}
or finally, as far as $\beta^5$, the equation is
\begin{align*}
0 &= +\beta(2\lambda - 20) + \beta^2(-19\lambda + 190) + \beta^3(153\lambda - 14\mu + 2\nu - 1138)\\
&\quad + \beta^4(-814\lambda + 119\mu - 17\nu + 4808) + \beta^5(3027\lambda - 558\mu + 94\nu - 15178).
\end{align*}

We have in the first place $\lambda = 10$, which makes the coefficient of $\beta^2$ to be $= 0$; and we then have
\[
7\mu - \nu - 196 = 0, \quad 119\mu - 17\nu - 3332 = 0, \quad 272\mu - 47\nu - 7546 = 0,
\]
where the second equation is equivalent to the first: the three equations give
\[
\mu = -\frac{833}{1}, \quad \nu = \frac{?}{1}.
\]

I have since found that proceeding one step further, that is, to $\beta^6$, the coefficient is
\[
-8310\lambda + 1791\mu - 363\nu + 36942,
\]
viz.\ putting this $= 0$, and for $\lambda$ substituting its value $= 10$, we have
\[
597\mu - 121\nu = 15386,
\]
an equation which is in fact satisfied by the values just obtained for $\lambda$ and $\mu$.

For the function $(Y, Z, W)$ we have the same equations with $\nu = 0$; and therefore $\lambda = 10$, $\mu = 28$; and for the function $(Y, Z)$ the same equations with $\mu = 0$, $\nu = 0$; and therefore $\lambda = 10$.

The linear combinations thus are
\begin{align*}
X &= a^{27} + bm^2,\\
Y &= a^{25}(c + b^2) + dm^2,\\
Y + 10Z &= a^{23}(e + c^2) + fm^2,\\
Y + 10Z + 28W &= a^{21}(g + d^2) + hm^2,\\
25Y + 250Z + 833W + 9312T &= a^{19}(i + e^2) + jm^2,
\end{align*}
where remark that in $Y + 10Z$, by means of the coefficient $10$, we make to disappear the two terms $a^{25}c$, $a^{25}d$; in the next function, by means of the coefficients $10$ and $28$, the four terms $a^{23}c$, $a^{23}d$, $a^{23}e$, $a^{23}f$, and in the last function, by means of the three coefficients, the six terms $a^{21}c$, $a^{21}d$, $a^{21}e$, $a^{21}f$, $a^{21}g$ and $a^{21}h$.

It is possible that a larger number of terms will disappear; but if this is so, the general form shown by the combinations on p.~353 will require modification.

The investigation applies without alteration to any odd weight $\omega$ whatever; the condition for the determination of the coefficients $\lambda$, $\mu$, $\nu$, \ldots\ is obviously $1 - 2\beta = 0$.

For the case of an even weight $\omega$, the series of functions is
\[
X = (\alpha - \beta)^{\omega}, \quad Y = (\alpha - \beta)^{\omega-2}(\alpha - \gamma)(\beta - \gamma), \quad Z = (\alpha - \beta)^{\omega-4}(\alpha - \gamma)^2(\beta - \gamma)^2, \quad \ldots
\]
and the condition for the determination of the coefficients $\lambda$, $\mu$, $\nu$, \ldots\ is in like manner found to be
\[
-1 + \lambda(0 - \beta^2) - \mu(0 - \beta^3) + \nu(0 - \beta^4) + \ldots = 0,
\]
which is to be understood as explained above.

\newpage

%==============================================================================
\section{Paper 848}
%==============================================================================

\begin{center}
\large\textbf{848.}

\Large\textbf{ON THE TRANSFORMATION OF THE DOUBLE-THETA FUNCTIONS.}

\smallskip
[From the \textit{Quarterly Journal of Pure and Applied Mathematics}, vol.~xxi. (1886), pp.~142--178.]
\end{center}

\bigskip

I propose to reproduce Hermite's Memoir ``Sur la th\'{e}orie de la transformation des fonctions Ab\'{e}liennes,'' Comptes Rendus, t.~XL. (1855), pp.~249, \ldots, 784, with some changes of notation and developments. Hermite's functions are even or odd according as we have $\mu q + \nu r$ even or odd: viz.\ his characteristic is $\binom{\mu, \nu}{q, r}$, or the letters $p$, $q$ are misplaced; I write, therefore, $r$ instead of $p$, so as to have the characteristic $\binom{\mu, \nu}{q, r}$; and then for symmetry it is necessary to interchange the suffixes $2$, $3$ and the letters $c$, $d$; the invariant function of the periods, instead of being as with him $\complement_0\complement_3(\complement_3\complement_0 + \complement_2\complement_1)\omega_1\nu_1$, must be taken to be $\complement_0\complement_2(\complement_2\complement_0 + \complement_1\complement_3)\omega_1\nu_1$.

Moreover, I write $A$, $B$ for his $G$, $\complement$, so as, instead of $(G, H, \complement \mid x, y)^2$, to have in the expressions of the theta-functions the quadric function $(A, H, B \mid x, y)^2$; and I alter the arrangement of the memoir so as to separate more completely the preliminary theory from the theory of the transformation.

\subsection*{General Theory. Art.~Nos. 1 to 21 (several sub-headings).}

\subsubsection*{The functions $\Pi(x, y \mid K, \text{indef.\ or def.})$.}

\textbf{Art.~Nos.~1--4.]
}

1.\ Consider a function $\Pi(x, y)(A, H, B \mid K, \text{indef.\ or def.})$, having: the characteristic $\binom{\mu, \nu}{q, r}$, where the characters $\mu$, $\nu$, $q$, $r$ are positive or negative integers, which may be taken to have each of them only the values $0$, $1$ at pleasure, so that the number of functions is $2^4 = 16$; the arguments $(x, y)$; the parameters, or conjoint quarter-periods, $(A, H, B)$; and the potency $K$, a positive integer; and which is either indefinite or definite, as will be explained; the function, moreover, contains linearly certain arbitrary constants, the number of them depending on the value of $K$, as will be explained.

The function may be written $\Pi(x, y)$ or in any other less abbreviated form which may be convenient.

\medskip

2.\ The function $\Pi(x, y \mid K \text{ indef.})$ is defined by the following four equations:
\begin{align*}
\Pi(x+1, y) &= (-1)^{\mu q}\Pi(x, y),\\
\Pi(x, y+1) &= (-1)^{\nu r}\Pi(x, y),\\
\Pi(x+A, y+H) &= (-1)^{q}\Pi(x, y)\exp(-2i\pi K(2x + A)),\\
\Pi(x+H, y+B) &= (-1)^{r}\Pi(x, y)\exp(-2i\pi K(2y + B));
\end{align*}
and the function $\Pi(x, y \mid K \text{ def.})$ by the same equations, together with the following fifth equation,
\[
\Pi(-x, -y) = (-1)^{\mu q + \nu r}\Pi(x, y);
\]
viz.\ the definite function is an even function or else an odd function of the arguments according as $\mu q + \nu r$ is even or odd.

We may call $\mu q + \nu r$ the \textit{index}; and the function is then even or odd according as the index is even or odd.

It is perhaps worth noticing that it would be allowable to define a function $\Pi(x, y \mid K, \text{skew def.})$, by the corresponding relation
\[
\Pi(-x, -y) = -(-1)^{\mu q + \nu r}\Pi(x, y),
\]
but I do not propose here to develope this notion.

\medskip

3.\ The four equations give rise to the following one,
\[
\Pi(x + a_0 + Aa_2 + Ha_3, y + a_1 + Ha_2 + Ba_3) = (-1)^{\mu a_0 + \nu a_1 + qa_2 + ra_3}\Pi(x, y) \times \exp(-2\pi K\{2a_2x + 2a_3y + (A, H, B \mid a_2, a_3)^2\}),
\]
where $a_0$, $a_1$, $a_2$, $a_3$ are any positive or negative integers (zero not excluded), and which single equation, in fact, includes the preceding four equations.

\medskip

4.\ In regard to the parameters it is to be observed that, if $A, H, B = A + i\alpha$, $H + i\eta$, $B + i\beta$, we must have $(\alpha, \eta, \beta \mid x, y)^2$ a determinate positive quadratic form; viz.\ this is the necessary and sufficient condition for the convergence of the series for the development of the function.

\subsubsection*{Number of arbitrary constants.}

\textbf{Art.~No.~5.}

5.\ The number of arbitrary constants is for the indefinite function $= K^2$; but for the definite function it is, when $K$ is odd, $= \frac{1}{2}(K^2 + 1)$; and when $K$ is even, it is $= \frac{1}{2}K^2$, except in the case of a characteristic $\binom{\mu, \nu}{q, r}$, when it is $= \frac{1}{4}(K^2 + 4)$.

In particular, for $K = 1$, there is only a single arbitrary constant, which is a mere factor of the function; taking it to be $= 1$, as presently explained, we have the 16 theta-functions.

\subsubsection*{Development in series.}

\textbf{Art.~Nos.~6--9.}

6.\ The function $\Pi(x, y)$ is developed in a series of exponentials, in the form
\[
\Pi(x, y) = \sum (-1)^{\mu q + \nu r}A_{m,n}\exp(i\pi\{(2m + \mu)x + (2n + \nu)y + (A, H, B \mid m, n)^2\}),
\]
where $m$ and $n$ have each of them all positive and negative integer values (zero not excluded) from $-\infty$ to $\infty$.

In fact, substituting this series in the four equations, they are all of them satisfied if only
\[
A_{m+K,\,n} = A_{m,\,n}, \quad A_{m,\,n+K} = A_{m,\,n}.
\]
Consequently the following $K^2$ coefficients remain arbitrary, viz.\ those with the suffixes $0$, $1$, \ldots, $K-1$; and we have for $\Pi(x, y)$ a sum of $K^2$ terms, each a determinate series multiplied into one of the arbitrary coefficients $A_{0,0}$, $A_{0,1}$, \&c.

The indefinite function thus contains, as already mentioned, $K^2$ arbitrary constants.

\medskip

7.\ Substituting in the fifth equation, we have for the definite function the further condition
\[
A_{-m,\,-n} = (-1)^{\mu q + \nu r}A_{m,\,n},
\]
which it is clear will be satisfied generally if only it is satisfied by the coefficients in the foregoing set of $K^2$ coefficients.

\medskip

8.\ In the case $K$ odd, we thus reduce the number of arbitrary coefficients to $\frac{1}{2}(K^2 + 1)$; the mode in which this takes place is best seen by an example. Suppose $K = 3$, so that $A_{m+3,\,n} = A_{m,\,n}$; $A_{m,\,n+3} = A_{m,\,n}$.

For the coefficients of the indefinite function, the suffixes are $00$, $01$, $02$, $10$, $11$, $12$, $20$, $21$, $22$.

And if we suppose $\mu = 0$, $\nu = 1$, then the new condition is $A_{-m,\,-n-1} = A_{m,\,n}$, viz.\ writing down only the suffixes, we thus obtain
\begin{align*}
00 &= 22, & 01 &= 21, & 02 &= 20,\\
10 &= 12, & 11 &= 11, & 12 &= 10,\\
20 &= 02, & 21 &= 01, & 22 &= 00,
\end{align*}
viz.\ one of these equations $01 = 01$ is an identity, but the other equations occur each twice; or we have four equations, each of them an equality between two out of the remaining 8 coefficients; the number of arbitrary coefficients is thus $1 + (9 - 1) = 5$; and so in general the number is $\frac{1}{2}(K^2 + 1)$.

\medskip

9.\ In the case $K$ even, the number is $= \frac{1}{2}K^2$, except for the characteristic $\binom{\mu, \nu}{q, r} = \binom{1, 1}{1, 1}$, when it is $= \frac{1}{4}(K^2 + 4)$.

\subsubsection*{The theta-functions.}

\textbf{Art.~Nos.~10--15.]
}

10.\ Writing $K = 1$, and taking $A_{0,0} = 1$, the value of the function thus is
\[
\vartheta_{\mu,\nu}(x, y) = \sum (-1)^{\mu q + \nu r}\exp(i\pi\{(2m + \mu)x + (2n + \nu)y + (A, H, B \mid m, n)^2\}).
\]

\medskip

11.\ The sum of two characteristics is the characteristic obtained by taking the sums of the component terms or characters,
\[
\binom{\mu}{q, r} + \binom{\mu'}{q', r'} = \binom{\mu + \mu'}{q + q', r + r'},
\]
and similarly for any number of characteristics.

I use the sign $=$, but this properly denotes a congruence, mod.~2; and the like as regards the indices.

The sum of two identical characteristics, or generally of any number of characteristics taken each of them any even number of times, is $= \binom{0}{0, 0}$. And this characteristic $\binom{0}{0, 0}$ may be called the characteristic $0$.

It should be observed, that the index of the sum is not in general equal to the sum of the indices. To make it so, we must have, for two characteristics,
\[
(\mu + \mu')(q + q') + (\nu + \nu')(r + r') = \mu q + \nu r + \mu' q' + \nu' r',
\]
that is,
\[
\mu q' + \mu' q + \nu r' + \nu' r = 0,
\]
and there is obviously a like formula for the case of more than two characteristics.

Two or more characteristics, such that they have the sum of the indices equal to the index of the sum, are said to be ``in direct relation'' or ``directly related'' to each other.

The sum of the indices and the index of the sum may differ by unity and we then have the inverse relation; but I do not propose to consider this.

\medskip

12.\ Consider any number $K$ of theta-functions, of the same arguments and parameters, but with the same or different characteristics. The product of these functions is in general a function $\Pi(x, y \mid K \text{ indef.})$, having a characteristic which is $=$ the sum of the characteristics of the theta-functions.

In fact, from the four equations of the theta-functions, we at once obtain for their products $\Pi(x, y)$ the four equations
\begin{align*}
\Pi(x+1, y) &= (-1)^{\Sigma\mu q}\Pi(x, y),\\
\Pi(x, y+1) &= (-1)^{\Sigma\nu r}\Pi(x, y),\\
\Pi(x+A, y+H) &= (-1)^{\Sigma q}\Pi(x, y)\exp(-i\pi K(2x + A)),\\
\Pi(x+H, y+B) &= (-1)^{\Sigma r}\Pi(x, y)\exp(-i\pi K(2y + B)),
\end{align*}
which proves the theorem.

\medskip

13.\ But if the indices are in direct relation to each other, then we have further
\[
\Pi(-x, -y) = (-1)^{\Sigma(\mu q + \nu r)}\Pi(x, y),
\]
and the product is thus a function $\Pi(x, y \mid K, \text{def.})$.

\medskip

14.\ Take the square of a theta-function, the characteristic is $= \binom{0}{0, 0}$ or $0$, and we have also twice the index $= 0$; viz.\ the theta-function is in direct relation with itself.

Hence the squared function is a function $\Pi(x, y \mid 2, \text{def.})$, and as such it contains linearly $\frac{1}{4}(2^2 + 4) = 4$ arbitrary constants.

Hence, taking any five squares, since each of them is a function of the form in question, it follows that the squares of the 5 theta-functions are connected by a linear relation.

\medskip

15.\ We may in a variety of ways (in fact, in 60 ways, as will presently be shown) select four theta-functions, all of them even, or else two of them even and two odd (that is, having the sum of their indices $= 0$), such that the sum of their characteristics is $= \binom{0}{0, 0}$; for instance, the functions may be
\[
\vartheta\binom{0,0}{0,0}, \quad \vartheta\binom{0,0}{0,1}, \quad \vartheta\binom{1,1}{0,0}, \quad \vartheta\binom{1,1}{0,1},
\]
indices $0$, $0$, $1$, $1$.

The functions are thus in direct relation, and the product of the four functions is a function $\Pi\binom{0}{0, 0}(4, \text{def.})$.

But obviously any one of the functions taken four times, or any two of them taken each twice, are in like manner four functions in direct relation, or the fourth powers $\vartheta^4$, $\vartheta'^4$, $\vartheta''^4$, $\vartheta'''^4$, and the squared products $\vartheta^2\vartheta'^2$, $\vartheta''^2\vartheta'''^2$, $\vartheta^2\vartheta''^2$, $\vartheta^2\vartheta'''^2$, $\vartheta'^2\vartheta''^2$, $\vartheta'^2\vartheta'''^2$, are in like manner each of them a function $\Pi\binom{0}{0, 0}(4, \text{def.})$, viz.\ we have thus in all $1 + 4 + 6 = 11$ such functions.

But the $\Pi$ function contains only $\frac{1}{4}(4^2 + 4) = 10$ arbitrary constants; hence there must be a linear relation between the 11 powers and products, and this is G\'{o}pel's relation.

\subsubsection*{The four functions $\Pi_0$, $\Pi_1$, $\Pi_2$, $\Pi_3$.}

\textbf{Art.~Nos.~16--21.]
}

16.\ Starting with any two characteristics $\mathbf{a}$, $\mathbf{b}$ at pleasure, the remaining characteristics form seven pairs, such that $\mathbf{a} + \mathbf{b} + \mathbf{c} + \mathbf{d} = 0$; but among the seven pairs we have only three, suppose $\mathbf{c} + \mathbf{d}$, $\mathbf{e} + \mathbf{f}$, $\mathbf{g} + \mathbf{h}$, which are such that $(\mathbf{a}, \mathbf{b}, \mathbf{c}, \mathbf{d})$, $(\mathbf{a}, \mathbf{b}, \mathbf{e}, \mathbf{f})$, $(\mathbf{a}, \mathbf{b}, \mathbf{g}, \mathbf{h})$ are each of them either all even or else two even and two odd; that is, starting with any pair $(\mathbf{a}, \mathbf{b})$, we have these three tetrads having each of them the required property.

The number of pairs $(\mathbf{a}, \mathbf{b})$ is $16 \cdot 15 = 240$; and we thence derive $240 \times 3 = 720$ tetrads; but each such tetrad is of course derivable from any one of the six pairs contained in it; or the number of distinct tetrads is $720/6 = 60$, viz.\ we have, as mentioned above, 60 G\'{o}pel-tetrads.

\medskip

17.\ We consider four theta-functions $\vartheta$, $\vartheta_1$, $\vartheta_2$, $\vartheta_3$, which are such that to the modulus 2, the sum of the characters is $= 0$, and also the sum of the indices is $= 0$; taking the characters to be
\[
\binom{\mu, \nu}{q, r}, \quad \binom{\mu', \nu'}{q', r'}, \quad \binom{\mu'', \nu''}{q'', r''}, \quad \binom{\mu''', \nu'''}{q''', r'''}
\]
and writing throughout $=$ for $=$ (mod.~2), we have
\begin{align*}
\mu + \mu' + \mu'' + \mu''' &= 0, & \nu + \nu' + \nu'' + \nu''' &= 0,\\
q + q' + q'' + q''' &= 0, & r + r' + r'' + r''' &= 0,\\
\mu q + \nu r + \mu' q' + \nu' r' + \mu'' q'' + \nu'' r'' + \mu''' q''' + \nu''' r''' &= 0.
\end{align*}

Writing for shortness
\[
(01) = \mu q' + \mu' q + \nu r' + \nu' r,
\]
and so in other cases; and further
\[
(01) + (02) + (12) = (012), \quad \text{\&c.,}
\]
then substituting for $\mu'''$, $\nu'''$, $q'''$, $r'''$ their values from the first four equations, we deduce
\[
(012) = 0; \quad \text{and similarly} \quad (013) = 0, \quad (023) = 0, \quad (123) = 0.
\]

\medskip

18.\ Consider now a product $\vartheta^a\vartheta_1^b\vartheta_2^c\vartheta_3^d$, where $a + b + c + d$ is $=$ a given odd number $k$; the characteristic is
\[
\binom{\mu a + \mu' b + \mu'' c + \mu''' d,\; \nu a + \nu' b + \nu'' c + \nu''' d}{qa + q'b + q''c + q'''d,\; ra + r'b + r''c + r'''d},
\]
and it hence follows that the index is $= a(\mu q + \nu r) + b(\mu' q' + \nu' r') + c(\mu'' q'' + \nu'' r'') + d(\mu''' q''' + \nu''' r''')$.

In fact, forming the index in question, we have first terms in $a^2$, $b^2$, $c^2$, $d^2$, which upon writing therein $a$, $b$, $c$, $d$ for these values respectively ($a^2 = a$, \&c.) give the required value; we have therefore only to show that the sum of the remaining terms in $ab$, \&c., is $= 0$.

These terms are
\[
ab(01) + ac(02) + ad(03) + bc(12) + bd(13) + cd(23),
\]
and writing herein $a + b + c + d = 1$, and thence $ad = a(1 - a - b - c) = ab + ac$, and similarly $bd = ab + bc$, $cd = ac + bc$, the terms in question become
\[
= ab(013) + ac(023) + bc(123),
\]
that is, they become $= 0$.

\medskip

19.\ We thus see that the function $\vartheta^a\vartheta_1^b\vartheta_2^c\vartheta_3^d$ has an index which is $= a\,\text{ind}\,\vartheta + b\,\text{ind}\,\vartheta_1 + c\,\text{ind}\,\vartheta_2 + d\,\text{ind}\,\vartheta_3$.

Consider separately four products $\vartheta^a\vartheta_1^b\vartheta_2^c\vartheta_3^d$, in which the exponents $a$, $b$, $c$, $d$ satisfy successively the relations (always to modulus 2)
\begin{align*}
b + d &= 0, & c + d &= 0,\\
b + d &= 1, & c + d &= 0,\\
b + d &= 0, & c + d &= 1,\\
b + d &= 1, & c + d &= 1.
\end{align*}
Combining herewith the relation $a + b + c + d = 1$, it follows that the exponents $a$, $b$, $c$, $d$ are
\begin{align*}
&= d + 1, d, d, d; & &= d, d + 1, d, d;\\
&= d, d, d + 1, d; & &= d, d, d, d + 1,
\end{align*}
in the four cases respectively.

Then substituting these values, the characteristics become
\[
\binom{\mu}{q, r}, \quad \binom{\mu'}{q', r'}, \quad \binom{\mu''}{q'', r''}, \quad \binom{\mu'''}{q''', r'''}
\]
viz.\ the four products have the same characters as $\vartheta$, $\vartheta_1$, $\vartheta_2$, $\vartheta_3$ respectively; and in like manner recollecting that $\text{ind}\,\vartheta + \text{ind}\,\vartheta_1 + \text{ind}\,\vartheta_2 + \text{ind}\,\vartheta_3 = 0$, we see that the four products have the same indices as $\vartheta$, $\vartheta_1$, $\vartheta_2$, $\vartheta_3$ respectively.

More generally write
\[
\Pi_0, \Pi_1, \Pi_2, \Pi_3 = \vartheta^a\vartheta_1^b\vartheta_2^c\vartheta_3^d,
\]
where for the four cases respectively the exponents $a$, $b$, $c$, $d$ satisfy the conditions already referred to; then $\Pi_0$, $\Pi_1$, $\Pi_2$, $\Pi_3$ have the same characteristics, and the same indices, as $\vartheta$, $\vartheta_1$, $\vartheta_2$, $\vartheta_3$ respectively.

\medskip

20.\ It can be shown that each of the functions $\Pi$ contains $\frac{1}{4}(k^2 + 1)$ constants.

It will be recollected that we have between $\vartheta$, $\vartheta_1$, $\vartheta_2$, $\vartheta_3$ an equation of the form
\[
a\vartheta^4 + b\vartheta_1^4 + c\vartheta_2^4 + d\vartheta_3^4 + e\vartheta^2\vartheta_1^2 + f\vartheta^2\vartheta_2^2 + \ldots = 0;
\]
this serves to express, say $\vartheta_3^4$, in lower powers of $\vartheta_3$; and by successive applications of this equation, we can reduce $\vartheta^a\vartheta_1^b\vartheta_2^c\vartheta_3^d$ to a form in which $d$ has only one of the values $0$, $1$, $2$ or $3$; we do not by this transformation alter the suffix of $\Pi$, viz.\ a term originally of the form $\Pi_0$, $\Pi_1$, $\Pi_2$ or $\Pi_3$, will by the transformation give rise only to terms which are of the same form $\Pi_0$, $\Pi_1$, $\Pi_2$ or $\Pi_3$ (as the case may be).

The number of constants in $\Pi_0$ is thus $=$ number of partitions of $k$ into four parts $a$, $b$, $c$, $d$, under the conditions
\begin{align*}
d &= 0 \text{ or } 2; & a \text{ odd}, b, c \text{ each even},\\
d &= 1 \text{ or } 3; & a \text{ even}, b, c \text{ each odd},
\end{align*}
where, in reckoning the partitions, the order of the parts is taken into account:
the partitions are thus as follows
\begin{align*}
d = 0, &\quad (a - 1) + b + c = k - 1,\\
d = 1, &\quad a + (b - 1) + (c - 1) = k - 3,\\
d = 2, &\quad (a - 1) + b + c = k - 3,\\
d = 3, &\quad a + (b - 1) + (c - 1) = k - 5,
\end{align*}
where the parts $a$ or $(a - 1)$, $b$ or $(b - 1)$, $c$ or $(c - 1)$, as the case may be, are all of them even; hence, writing $k' = \frac{1}{2}(k - 1)$, the cases are
\[
a' + b' + c' = k', \quad k' - 1, \quad k' - 1, \quad \text{or } k' - 2,
\]
where the $a'$, $b'$, $c'$ are odd or even (zero not excluded) at pleasure; as already mentioned, the order of the parts is taken into account: thus the particulars of $3$ would be
\begin{align*}
&300, \quad 210, \quad 120, \quad 030,\\
&201, \quad 111, \quad 021, \quad 102, \quad 012, \quad 003.
\end{align*}
No.\ is $10 = 4 \cdot 5$.

Hence, in the four cases respectively, the numbers are
\[
\tfrac{1}{2}(k'^2 + k' + 2), \quad \tfrac{1}{2}(k'^2 + k'), \quad \tfrac{1}{2}(k'^2 + k'), \quad \tfrac{1}{2}(k'^2 - k'),
\]
giving a total
\[
= \tfrac{1}{2}(4k'^2 + 4k' + 2) = \tfrac{1}{2}\{(2k' + 1)^2 + 1\},
\]
that is, $= \frac{1}{2}(k^2 + 1)$.

And similarly the number is $= \frac{1}{2}(k^2 + 1)$ in the other three cases respectively.

\medskip

21.\ It can be shown that, conversely, from any function $\Pi(x, y \mid K, \text{def.})$ having the characteristic and index of $\vartheta$, we can obtain the several functions $\Pi_0$, $\Pi_1$, $\Pi_2$, $\Pi_3$ as linear functions of the products $\vartheta^a\vartheta_1^b\vartheta_2^c\vartheta_3^d$; or, what is the same thing, the function $\Pi(x, y \mid K, \text{def.})$ is expressible as a linear function of the products in question.

\newpage

\subsection*{Preparation for the Transformation. Art.~Nos.\ 22 to 44 (several sub-headings).}

\subsubsection*{The Hermitian quartic matrix.}

\textbf{Art.~Nos.~22--25.]
}

22.\ Observing that, for the adopted form of the $\Pi$- or $\vartheta$-functions, the periods $\omega$, $\nu$ are
\[
\begin{matrix}
1, & 0, & \omega_1, & \nu_1\\
0, & 1, & \omega_2, & \nu_2\\
A, & H, & \omega_3, & \nu_3
\end{matrix}
\]
so that we have
\[
\omega_1\nu_2 - \omega_2\nu_1 = A, \quad \omega_1\nu_3 - \omega_3\nu_1 = H, \quad \omega_2\nu_3 - \omega_3\nu_2 = B,
\]
we have to consider the automorphic transformation of the bilinear form $\omega_1\nu_2 - \omega_2\nu_1 + \omega_1\nu_3 - \omega_3\nu_1$.

\medskip

23.\ We write
\[
(\omega_0, \omega_1, \omega_2, \omega_3 \mid \nu_0, \nu_1, \nu_2, \nu_3) = (a_0, a_1, a_2, a_3 \mid b_0, b_1, b_2, b_3)(\tau_0, \tau_1, \tau_2, \tau_3),
\]
and the coefficients are assumed to be such that we have identically
\[
\omega_0\nu_1 - \omega_1\nu_0 + \omega_2\nu_3 - \omega_3\nu_2 = k(\tau_0\tau_1 - \tau_1\tau_0 + \tau_2\tau_3 - \tau_3\tau_2),
\]
where $k$ is in the sequel taken to be a positive integer.

We obtain by direct substitution the value of $\omega_0\nu_1 - \omega_1\nu_0 + \omega_2\nu_3 - \omega_3\nu_2$ in the following form:
\begin{align*}
&(a_2c_0 - c_2a_0 + b_2d_0 - d_2b_0)(\tau_0\tau_1 - \tau_1\tau_0)\\
&\quad + (a_3c_0 - c_3a_0 + b_3d_0 - d_3b_0)(\tau_0\tau_2 - \tau_2\tau_0)\\
&\quad + (a_0c_1 - c_0a_1 + b_0d_1 - d_0b_1)(\tau_0\tau_3 - \tau_3\tau_0)\\
&\quad + (a_3c_1 - c_3a_1 + b_3d_1 - d_3b_1)(\tau_1\tau_2 - \tau_2\tau_1)\\
&\quad + (a_1c_2 - c_1a_2 + b_1d_2 - d_1b_2)(\tau_1\tau_3 - \tau_3\tau_1)\\
&\quad + (a_0c_3 - c_0a_3 + b_0d_3 - d_0b_3)(\tau_2\tau_3 - \tau_3\tau_2);
\end{align*}
viz.\ equating this to its value $k(\tau_0\tau_1 - \tau_1\tau_0 + \tau_2\tau_3 - \tau_3\tau_2)$, we have 4 identities and 12 equations which are, in fact, 6 equations occurring each twice.

We have thus six equations, which are the conditions in order that the matrix may be the matrix of automorphic transformation of the bilinear form $\omega_0\nu_1 - \omega_1\nu_0 + \omega_2\nu_3 - \omega_3\nu_2$.

The six equations may be written
\begin{align*}
(ac + bd)_{01} &= 0, & (ac + bd)_{02} &= k, & (ac + bd)_{03} &= 0,\\
(ac + bd)_{12} &= 0, & (ac + bd)_{13} &= k, & (ac + bd)_{23} &= 0,
\end{align*}
viz.\ the first of these equations is $a_0c_1 - a_1c_0 + b_0d_1 - b_1d_0 = 0$, and so in other cases.

It is convenient to remark that each interchange of two letters $a$ and $c$, $b$ and $d$, also interchange of the suffixes $0$ and $2$, produces a change of sign; thus the second equation may be written $(ca + db)_{02} = k$, or $(ca + db)_{20} = k$.

\medskip

24.\ The inverse matrix is found to be
\[
\begin{pmatrix}
a_0, & a_1, & a_2, & a_3\\
b_0, & b_1, & b_2, & b_3\\
c_0, & c_1, & c_2, & c_3\\
d_0, & d_1, & d_2, & d_3
\end{pmatrix}^{-1}
= \frac{1}{k}
\begin{pmatrix}
c_1, & d_1, & -a_1, & -b_1\\
-c_0, & -d_0, & a_0, & b_0\\
c_3, & d_3, & -a_3, & -b_3\\
-c_2, & -d_2, & a_2, & b_2
\end{pmatrix},
\]
and the determinant of each of the matrices in this formula is $= k^2$.

We have thus
\[
(\nu_0, \nu_1, \nu_2, \nu_3) = \frac{1}{k}
\begin{pmatrix}
c_1, & d_1, & -a_1, & -b_1\\
-c_0, & -d_0, & a_0, & b_0\\
c_3, & d_3, & -a_3, & -b_3\\
-c_2, & -d_2, & a_2, & b_2
\end{pmatrix}
(\omega_0, \omega_1, \omega_2, \omega_3),
\]
and the like formula for the $\tau$, $\upsilon$.

Substituting these values in the equation
\[
k(\tau_2\upsilon_2 - \tau_2\upsilon + \tau_1\upsilon_2 - \tau_3\upsilon_1) = \omega_0\nu_1 - \omega_1\nu_0 + \omega_2\nu_3 - \omega_3\nu_2,
\]
we obtain 6 new equations, which are in a different form the conditions for the automorphic transformation.

The 6 new equations may be written
\begin{align*}
(ac + bd)_{01} &= 0, & (ac + bd)_{02} &= k, & (ac + bd)_{03} &= 0,\\
(ac + bd)_{12} &= 0, & (ac + bd)_{13} &= k, & (ac + bd)_{23} &= 0,
\end{align*}
viz.\ these equations are $(ac + bd)_{01} = 0$, $(ac + bd)_{02} = k$, $(ac + bd)_{03} = 0$, $(ac + bd)_{12} = 0$, $(ac + bd)_{13} = k$, $(ac + bd)_{23} = 0$, or say
\[
a_0c_1 - a_1c_0 + b_0d_1 - b_1d_0 = 0, \quad \text{\&c.}
\]

\medskip

25.\ It is worth while to show how the foregoing formula for the inverse matrix comes out.

Take for instance the diagonal minor
\[
\begin{vmatrix}
b_1, & b_2, & b_3\\
c_1, & c_2, & c_3\\
d_1, & d_2, & d_3
\end{vmatrix}
= b_1(c_2d_3 - c_3d_2) + b_2(c_3d_1 - c_1d_3) + b_3(c_1d_2 - c_2d_1),
\]
which is $= c_1(ac)_{23} + c_2(ac)_{31} + c_3(ac)_{12}$, since the remaining terms destroy each other. And dividing by the determinant, which is $= k^2$, we have the term $c_2/k$ of the inverse matrix.

\subsubsection*{The Symmetrical Hermitian Matrix.}

\textbf{Art.~Nos.~26--28.]
}

26.\ We may consider a symmetrical Hermitian matrix, say the matrix
\[
\begin{pmatrix}
\mathfrak{A}, & \mathfrak{B}, & \mathfrak{C}, & \mathfrak{D}\\
\mathfrak{B}, & \mathfrak{E}, & \mathfrak{F}, & \mathfrak{G}\\
\mathfrak{C}, & \mathfrak{F}, & \mathfrak{H}, & \mathfrak{I}\\
\mathfrak{D}, & \mathfrak{G}, & \mathfrak{I}, & \mathfrak{J}
\end{pmatrix},
\]
viz.\ we have
\begin{align*}
2\mathfrak{B} - \mathfrak{E} + \mathfrak{H} &= 0, & 2\mathfrak{C} - \mathfrak{F} + \mathfrak{I} &= 0,\\
2\mathfrak{D} - \mathfrak{G} + \mathfrak{J} &= 0, & 2\mathfrak{F} - \mathfrak{G} + \mathfrak{H} &= 0,\\
2\mathfrak{C} - \mathfrak{D} + \mathfrak{I} &= 0, & 2\mathfrak{B} - \mathfrak{C} + \mathfrak{F} &= 0.
\end{align*}
The characteristic property is that, effecting a Hermitian transformation, we have a new symmetrical Hermitian matrix.

In fact, the new matrix is
\[
\begin{pmatrix}
a_0, & b_0, & c_0, & d_0\\
a_1, & b_1, & c_1, & d_1\\
a_2, & b_2, & c_2, & d_2\\
a_3, & b_3, & c_3, & d_3
\end{pmatrix}
\begin{pmatrix}
\mathfrak{A}, & \mathfrak{B}, & \mathfrak{C}, & \mathfrak{D}\\
\mathfrak{B}, & \mathfrak{E}, & \mathfrak{F}, & \mathfrak{G}\\
\mathfrak{C}, & \mathfrak{F}, & \mathfrak{H}, & \mathfrak{I}\\
\mathfrak{D}, & \mathfrak{G}, & \mathfrak{I}, & \mathfrak{J}
\end{pmatrix}
\begin{pmatrix}
a_0, & a_1, & a_2, & a_3\\
b_0, & b_1, & b_2, & b_3\\
c_0, & c_1, & c_2, & c_3\\
d_0, & d_1, & d_2, & d_3
\end{pmatrix},
\]
where the first factor, qua transposed matrix, is Hermitian. Hence the product is a symmetrical $\frac{1}{k}$-matrix.

\medskip

27.\ Write
\[
\Delta = \alpha\beta - \eta^2 = (\alpha, \eta, \beta \mid x, y)^2,
\]
and consider the following matrix
\[
\begin{pmatrix}
\eta, & \alpha, & \eta, & \alpha\\
\beta, & \eta, & \beta, & \eta\\
\eta, & \alpha, & \eta, & \alpha\\
\beta, & \eta, & \beta, & \eta
\end{pmatrix}
= (p, q, r, s \mid x, y, z, w),
\]
where
\begin{align*}
p &= \eta x + \alpha y + \eta z + \alpha w,\\
q &= \beta x + \eta y + \beta z + \eta w,\\
r &= \eta x + \alpha y + \eta z + \alpha w,\\
s &= \beta x + \eta y + \beta z + \eta w.
\end{align*}

It is easily shown that this is a symmetrical $\Delta$-matrix. In fact, representing it for a moment by
\[
\begin{pmatrix}
\mathfrak{A}, & \mathfrak{B}, & \mathfrak{C}, & \mathfrak{D}\\
\mathfrak{B}, & \mathfrak{E}, & \mathfrak{F}, & \mathfrak{G}\\
\mathfrak{C}, & \mathfrak{F}, & \mathfrak{H}, & \mathfrak{I}\\
\mathfrak{D}, & \mathfrak{G}, & \mathfrak{I}, & \mathfrak{J}
\end{pmatrix},
\]
so that $\mathfrak{A} = \eta$, $\mathfrak{E} = \eta$, $\mathfrak{H} = \eta$, \&c., we have
\[
2\mathfrak{B} - \mathfrak{E} + \mathfrak{H} = (\Delta - \beta\mathfrak{A}_0) - (-\eta^2 - \eta\mathfrak{A}_0) - \eta^2\mathfrak{A} - \beta\Delta + \eta\mathfrak{B}_0 + \eta^2\mathfrak{A}_0,
\]
which is $= \alpha\beta\Delta - \beta\delta\mathfrak{A}_0 + \eta^2\mathfrak{A}_0$, that is, $= (\alpha\beta - \eta^2)\Delta = \Delta^2$.

And, similarly, $\mathfrak{A}\mathfrak{H} - \mathfrak{C}^2 + \mathfrak{E}\mathfrak{I} - \mathfrak{F}^2$ is found to be $= \Delta^2$; and the remaining four combinations of terms to be each of them $= 0$; the matrix is thus a symmetrical $\Delta$-matrix.

\medskip

28.\ It is to be added that the diagonal minors and the determinant
\begin{align*}
\mathfrak{A}, &\quad \mathfrak{A}\mathfrak{E} - \mathfrak{B}^2,\\
\mathfrak{A}\mathfrak{E}\mathfrak{H} + \text{\&c.}, &\quad \mathfrak{A}\mathfrak{E}\mathfrak{H}\mathfrak{J} + \text{\&c.},
\end{align*}
have respectively the values $\Delta$, $\Delta^2$, $\Delta^3$, $\Delta^4$; viz.\ if $\alpha$, $\beta$, $\delta$ are positive, then these are all positive; or the last-mentioned quadric function is a definite positive form.

\newpage

\subsubsection*{The general matrix resumed: Arithmetical theory.}

\textbf{Art.~Nos.~29--36.]
}

29.\ The matrix containing, as before, the parameter $k$, may be called a $k$-matrix; if $k$ be $= 1$, it is a unit-matrix.

From the fundamental equation
\[
k(\tau_0\upsilon_1 - \tau_1\upsilon_0 + \tau_2\upsilon_3 - \tau_3\upsilon_2) = \omega_0\nu_1 - \omega_1\nu_0 + \omega_2\nu_3 - \omega_3\nu_2,
\]
it at once follows that, compounding a $k$-matrix with a $k'$-matrix, we have a $kk'$-matrix; and in particular, compounding a $k$-matrix with a unit-matrix, we have a $k$-matrix.

\medskip

30.\ The symbols $a_0$, $a_1$, $b_0$, \&c., and $k$, have thus far been arbitrary magnitudes, but we now take them to be integers; and we consider in particular the case where $k$ is a positive odd prime.

The number of $k$-matrices is of course infinite, but if we regard as equivalent any two such matrices which are derivable one from the other by post-multiplication by a unit-matrix (viz.\ $U$ being a unit-matrix, the matrices $M$ and $M \cdot U$ are regarded as equivalent), then the number of distinct $k$-matrices is finite, and $= 1 + k + k^2 + k^3$.

The first step is to show that we can by post-multiplication, by a properly determined unit-matrix, reduce the $k$-matrix to the form
\[
\begin{pmatrix}
a_0, & a_1, & a_2, & a_3\\
b_0, & b_1, & b_2, & b_3\\
c_0, & 0, & c_2, & 0\\
d_0, & 0, & 0, & d_3
\end{pmatrix},
\]
these values being such as to satisfy identically two out of the six conditions; the remaining conditions present themselves under the two equivalent forms
\begin{align*}
a_0c_2 &= k, & b_1d_3 &= k,\\
a_0b_2 + a_1c_2 - a_2c_1 - a_3c_0 &= 0, & a_0d_1 + a_1d_0 - a_2d_3 - a_3d_2 &= 0,
\end{align*}
and
\begin{align*}
a_0c_2 &= k, & b_1d_3 &= k,\\
a_0c_2 + b_1c_2 - b_2c_1 - b_3c_0 &= 0, & a_3c_2 + b_2d_3 - b_3d_2 - b_0d_1 &= 0.
\end{align*}

Hence $a_0$, $c_2 = 1, k$, or $k, 1$; and $b_1$, $d_3 = 1, k$, or $k, 1$; so that, combining these pairs of values, we have four different types of matrix, each type depending on the coefficients $a_1$, $d_0$, $a_2$, $b_2$, $a_3$, $b_3$, connected together by two equations.

But the forms of the same type are not distinct from each other, and we have to determine for each of the four types a system of non-equivalent forms comprised therein, and such that from these, by post-multiplication by a unit-matrix as before, the other forms of the type can be obtained.

This final system is:

\medskip

\textbf{Type I.}
\[
\begin{pmatrix}
1, & 0, & 0, & 0\\
0, & 1, & 0, & 0\\
0, & 0, & k, & 0\\
0, & 0, & 0, & k
\end{pmatrix}
\]

\textbf{Type II.}
\[
\begin{pmatrix}
1, & 0, & 0, & 0\\
0, & k, & 0, & 0\\
0, & i, & 1, & 0\\
0, & 0, & 0, & k
\end{pmatrix}
\]

\textbf{Type III.}
\[
\begin{pmatrix}
k, & 0, & 0, & 0\\
i, & 1, & 0, & 0\\
i', & 0, & 1, & 0\\
0, & 0, & 0, & k
\end{pmatrix}
\]

\textbf{Type IV.}
\[
\begin{pmatrix}
k, & 0, & 0, & 0\\
i, & 1, & 0, & 0\\
i', & 0, & 1, & 0\\
i'', & 0, & 0, & 1
\end{pmatrix}
\]

\medskip

where $i$, $i'$, $i''$ are integers, having each of them any one of the values $0$, $1$, $2$, \ldots, $k-1$, viz.\ there are in the four types $1$, $k$, $k^2$, $k^3$ forms respectively, and the number of forms is thus $= 1 + k + k^2 + k^3$, as already mentioned.

I abstain from the further details of the proof.

There is obviously a like theory for which, in place of post-multiplication, we have pre-multiplication; viz.\ here the matrices $M$ and $U \cdot M$ are regarded as equivalent.

\medskip

31.\ Any two $k$-matrices are reducible one to the other by a combined pre- and post-multiplication; viz.\ we have always $M = U \cdot M' \cdot U'$, where $M$, $M'$ are any two $k$-matrices, and $U$, $U'$ properly determined unit-matrices; and in particular, $M$ being any given $k$-matrix, this is expressible in the foregoing form, where $M'$ denotes the principal matrix
\[
\begin{pmatrix}
1, & 0, & 0, & 0\\
0, & 1, & 0, & 0\\
0, & 0, & k, & 0\\
0, & 0, & 0, & k
\end{pmatrix}.
\]

\subsubsection*{Congruence theorems, $k$ an odd number.}

\medskip

32.\ Taking $k$ an odd number, and using throughout $=$ instead of $\equiv$ (mod.~2), we have the following congruences:
\begin{align*}
\mu &\equiv a(a_2c_2 + a_3c_3) + b(b_2d_2 + b_3d_3),\\
\nu &\equiv a(a_2c_0 + a_3c_1) + b(b_2d_0 + b_3d_1),\\
q &\equiv c(a_0c_2 + a_1c_3) + d(b_0d_2 + b_1d_3),\\
r &\equiv c(a_0c_0 + a_1c_1) + d(b_0d_0 + b_1d_1),
\end{align*}
or, conversely,
\begin{align*}
a &\equiv c(c_2\mu + c_3\nu) + d(d_2q + d_3r),\\
b &\equiv c(c_0\mu + c_1\nu) + d(d_0q + d_1r),\\
c &\equiv a(a_2\mu + a_3q) + b(b_2\nu + b_3r),\\
d &\equiv a(a_0\mu + a_1q) + b(b_0\nu + b_1r),
\end{align*}
where observe that on the right-hand side the signs may be changed into $+$; in fact, to the modulus 2, we have for any integer value whatever $p \equiv +p$.

\medskip

33.\ The first congruence is
\[
\mu \equiv a(a_2c_2 + a_3c_3) + b(b_2d_2 + b_3d_3), \quad \text{say } X = Y,
\]
and to verify this, I see no other method than that of considering separately all the combinations of even and odd values of $a$, $a_1$, $a_2$, $a_3$; viz.\ we have the 16 cases

\medskip
\begin{center}
\begin{tabular}{c|cccc|cccc|c}
Case & $a$ & $a_1$ & $a_2$ & $a_3$ & $b$ & $b_1$ & $b_2$ & $b_3$ & $Y$ \\
\hline
1 & e & e & e & e & e & e & e & e & 0 \\
2 & e & e & e & e & o & o & o & o & 0 \\
3 & o & o & o & o & e & e & e & e & $c_2 + c_3 + d_2 + d_3$ \\
4 & o & o & o & o & o & o & o & o & $c_2 + c_3 + d_2 + d_3 + b_1d_1 + b_3d_3$ \\
5 & e & e & o & o & e & e & e & e & 0 \\
6 & e & e & o & o & o & o & o & o & 0 \\
7 & o & o & e & e & e & e & e & e & $c_0 + c_1 + d_0 + d_1$ \\
8 & o & o & e & e & o & o & o & o & $c_0 + c_1 + d_0 + d_1 + b_1d_1 + b_3d_3$ \\
9 & e & o & e & o & e & o & e & o & 0 \\
10 & e & o & e & o & o & e & o & e & 0 \\
11 & o & e & o & e & e & o & e & o & $c_1 + c_3 + d_1 + d_3$ \\
12 & o & e & o & e & o & e & o & e & $c_1 + c_3 + d_1 + d_3 + b_0d_0 + b_2d_2$ \\
13 & e & o & o & e & e & o & o & e & 0 \\
14 & e & o & o & e & o & e & e & o & 0 \\
15 & o & e & e & o & e & o & o & e & $c_0 + c_2 + d_0 + d_2$ \\
16 & o & e & e & o & o & e & e & o & $c_0 + c_2 + d_0 + d_2 + b_0d_0 + b_2d_2$
\end{tabular}
\end{center}
\medskip

34.\ Case 2.\ We have $Y = b_1d_1$ which should be $= 0$, and in fact, the six equations give $b_1 = 0$, $d_1 = 0$, whence $Y = 0$.

Case 4.\ We have $Y = d_0 + c_0 + b_1d_1 + b_3d_3$ which should $= 1$, and in fact, the six equations give $c_3 - c_1 = 1$, $b_3 - b_1 = 0$, $d_3 - d_1 = 0$; whence $d_0 + c_0 = 1$, $b_1d_1 + b_3d_3 = 0$, and therefore $Y = 1$.

Case 6.\ Here $Y = b_0d_0 + b_1d_1$ which should $= 0$; and in fact, the six equations give $b_0 + b_1 = 0$, $d_0 + d_1 = 0$, whence $Y = 0$.

Case 16.\ Here $Y = c_0 + d_0 + c_2 + c_3 + b_0d_0 + b_1d_1 + b_2d_2 + b_3d_3$, which should be $= 0$; the six equations give
\[
-c_0 - d_0 + c_2 + c_3 = 1, \quad -b_0 - b_1 + b_2 + b_3 = 0, \quad -d_0 - d_1 + d_2 + d_3 = 0,
\]
and $b_2c_0 - b_0c_2 + b_3c_1 - b_1c_3 = 1$.

Writing $b_0 + b_2 = b_1 + b_3$ and $d_0 + d_2 = d_1 + d_3$, we find
\[
b_0d_0 + b_2c_0 + b_2d_2 + b_3d_3 = b_1d_1 + b_3d_3 + b_2d_0 + b_0d_2,
\]
that is,
\[
b_0d_0 + b_1d_1 + b_2d_2 + b_3d_3 = b_2d_0 + b_0d_2 + b_1d_1 + b_3d_3 = 1;
\]
and $c_0 + c_1 + c_2 + c_3 = 1$, whence $Y = 0$.

\medskip

35.\ Case 8.\ This is the only case of any difficulty: we have
\[
Y = c_0 + c_1 + d_0 + d_1 + b_1d_1 + b_3d_3,
\]
which should be $= 1$.

The six equations give
\begin{align*}
c_0 &= -1 - c_1 + c_3, & b_0 &= -b_1 + b_3,\\
d_0 &= -d_1 + d_3, & b_2 &= b_2, & c_2 &= c_2, & d_2 &= d_2,
\end{align*}
or, say
\[
c_0 = -1 - c_1 + c_3, \quad b_0 = -b_1 + b_3, \quad d_0 = -d_1 + d_3,
\]
or, substituting these values and omitting even terms
\[
Y = c_1 + c_3 - b_3d_3 - b_1d_1.
\]

The remaining three of the six equations are
\[
b_2c_0 - b_0c_2 + b_3c_1 - b_1c_3 = 1, \quad b_2d_0 - b_0d_2 + b_3d_1 - b_1d_3 = 0, \quad c_2d_0 - c_0d_2 + c_3d_1 - c_1d_3 = 0;
\]
or, substituting for $b_0$, $c_0$, $d_0$ their values, these become
\begin{align*}
(1 + c_1 - c_3)b_2 + (b_3 - b_1)c_2 &= -b_3 + b_2c_1,\\
(d_1 - d_3)b_2 + (b_3 - b_1)d_2 &= 1 - b_3 + b_1d_1,\\
(d_1 - d_3)c_2 + (-1 - c_1 + c_3)d_2 &= -d_3 + c_1d_1;
\end{align*}
we can from these equations eliminate $b_2$, $c_2$, $d_2$; viz.\ from the first and third equations eliminating $c_2$, we have
\[
(d_3 - d_1)\{b_2(1 + c_1 - c_3) + c_2(b_3 - b_1)\} = (d_3 - d_1)(-b_3 + b_2c_1) + (b_3 - b_1)(-d_3 + c_1d_1),
\]
and this, by means of the second equation, becomes
\[
(1 + d_1 - c_3)(-1 + b_3 - b_1) = (d_3 - d_1)(-b_3 + b_3c_1) + (b_3 - b_1)(-c_1d_1 + c_1d_1),
\]
viz.\ reducing, this is
\[
-d_1 + c_3 - 1 + b_3 - b_3d_1 = 0;
\]
and in virtue of it we have $Y = c_1 + c_3 - b_3d_3 - b_1d_1 = 1$.

It can be further shown that, to the modulus 2 ($k$ being, as before, odd), we have
\begin{align*}
&(a_0a_2 + a_1a_3)(c_0c_2 + c_1c_3) + (b_0b_2 + b_1b_3)(d_0d_2 + d_1d_3) = 0,\\
&(a_0c_0 + b_0d_0)(a_2c_2 + b_2d_2) + (a_1c_1 + b_1d_1)(a_3c_3 + b_3d_3) = 0.
\end{align*}
To prove the first equation, write for a moment
\[
X = (a_0c_1 - a_1c_0)(a_2c_3 - a_3c_2), \quad Y = (b_0d_1 - b_1d_0)(b_2d_3 - b_3d_2),
\]
then in virtue of the equations we have $X = X'$. But we have identically
\[
XY - X'Y' = (a_0c_2 + a_1c_3)(a_2c_0 + a_3c_1) + (b_0d_2 + b_1d_3)(b_2d_0 + b_3d_1),
\]
and from the equation $a_0c_2 - a_2c_0 = 1 + a_1c_3 - a_3c_1$, we have $XY = 0$; and similarly $X'Y' = 0$, that is, $XY = X'Y' = 0$; we have thus the required equation $XY + X'Y' = 0$.

In a similar manner the second equation may be verified.

\medskip

36.\ Write
\begin{align*}
p &= \mu a_0 + \nu a_1 + qa_2 + ra_3 + aa_2 + ba_3,\\
\nu' &= \mu b_0 + \nu b_1 + qb_2 + rb_3 + bb_2 + bb_3,\\
q' &= \mu c_0 + \nu c_1 + qc_2 + rc_3 + cc_2 + cc_3,\\
r' &= \mu d_0 + \nu d_1 + qd_2 + rd_3 + dd_2 + dd_3.
\end{align*}
It is to be shown that to the modulus 2 we have
\[
\mu q' + \nu r' = \mu q + \nu r.
\]

In fact, forming the value of $\mu q' + \nu r'$, we have first a constant term (term without $\mu$, $\nu$, $q$, $r$) which vanishes; next writing $\mu^2 = \mu$, the whole term in $\mu$ is
\[
c_0(a_0a_2 + a_1a_3) + d_0(b_0b_2 + b_1b_3) + c_2(a_0a_2 + a_1a_3) + d_2(b_0b_2 + b_1b_3),
\]
which also vanishes; and similarly the terms in $\nu$, $q$, $r$ each of them vanish; there remain only the terms in $\mu\nu$, $\mu q$, \&c.

The coefficient of $\mu q$ is $a_0c_2 + a_1c_0 + b_0d_2 + b_1d_0$, which (always to the modulus 2) is $= a_0c_2 - a_2c_0 + b_0d_2 - b_2d_0$, that is, it is $= 1$; and similarly the coefficient of $\nu r$ is $= 1$; and in like manner the coefficients of the other terms are each $= 0$; we have thus the required congruence
\[
\mu q' + \nu r' = \mu q + \nu r.
\]

\subsubsection*{The quintic matrix.}

\textbf{Art.~Nos.~37--43.]
}

37.\ I consider a quintic matrix composed of the foregoing coefficients $(a, b, c, d)_0^{1, 2, 3}$, viz.\ the matrix contained in a linear transformation which I write as follows:
\[
(T, P, Q, R, S) =
\begin{pmatrix}
(ab)_{01}, & (ab)_{21}, & 2(ab)_{31}, & (ab)_{13}, & (ab)_{32}\\
(ac)_{01}, & (ac)_{21}, & (ac)_{31} + k, & (ac)_{13}, & (ac)_{32}\\
(ad)_{01}, & (ad)_{21}, & 2(ad)_{31}, & (ad)_{13}, & (ad)_{32}\\
(cb)_{01}, & (cb)_{21}, & 2(cb)_{31}, & (cb)_{13}, & (cb)_{32}\\
dc_{01}, & dc_{21}, & 2dc_{31}, & dc_{13}, & dc_{32}
\end{pmatrix}
(T, P, Q, R, S),
\]
where as before $(ab)_{ij} = a_ib_j - a_jb_i$, \&c.; and so in other cases; in particular, observe that, in the expression of $Q$, the term involving $Q$ is $\{k + 2(ac)_{31}\}Q$ which, in virtue of the relation $(ac + bd)_{31} = k$, may also be written $\{(ac)_{31} - (bd)_{31}\}Q$.

I notice that in Hermite's paper, p.~366, the term is in effect written without the $k$, $= 2(ac)_{31}Q$: the correction of this erratum and of a corresponding one, p.~366, is made p.~787 at the conclusion of the memoir.

\medskip

38.\ The matrix is automorphic for the form $Q^2 - PR - TS$, viz.\ we have identically
\[
Q'^2 - P'R' - T'S' = k^2(Q^2 - PR - TS).
\]
As a partial verification, consider in $Q'^2 - P'R' - T'S'$ the term containing $Q^2$.

The coefficient of $Q^2$ is
\[
\{-k + 2(ac)_{02}\}^2 - 4(ab)_{02}(ad)_{02} - 4(ab)_{01}(dc)_{02},
\]
where the first term is
\[
k^2 - 4k(ac)_{02} + 4(ac)_{02}^2.
\]
Hence, observing that we have
\[
(ac)_{02}^2 - (ab)_{02}(ad)_{02} - (ab)_{01}(dc)_{02} = k(ac)_{02},
\]
the whole coefficient is $= k^2$, as it should be; and in like manner the verification may be effected for any other term.

\medskip

39.\ We require the following formulae:
\begin{align*}
(-c_0, a_0, b_0)(T, P, Q) &= k(b_1, -b_0, -b_3)(T, P, Q, R, S),\\
(-c_1, a_1, b_1)(T, P, Q) &= k(., b_0, -b_2)(T, P, Q, R, S),\\
(-c_2, a_2, b_2)(T, P, Q) &= k(b_3, ., b_0)(T, P, Q, R, S),\\
(-c_3, a_3, b_3)(T, P, Q) &= k(-b_2, b_1, .)(T, P, Q, R, S),
\end{align*}
that is,
\[
-c_0T + a_0P + b_0Q = k(b_1P - b_0Q - b_3S), \quad \text{\&c.,}
\]
and
\begin{align*}
(-d_0, a_0, b_0)(T, Q, R) &= k(., -a_0, ., a_3)(T, P, Q, R, S),\\
(-d_1, a_1, b_1)(T, Q, R) &= k(., a_1, ., -a_2)(T, P, Q, R, S),\\
(-d_2, a_2, b_2)(T, Q, R) &= k(., -a_3, ., a_0)(T, P, Q, R, S),\\
(-d_3, a_3, b_3)(T, Q, R) &= k(., a_2, ., -a_1)(T, P, Q, R, S).
\end{align*}

From the first set, multiplying the first, third and fourth equations by $T$, $P$, $Q$ respectively and adding, and again multiplying the second, third and fourth equations by $T$, $Q$, $R$, and adding, we obtain
\begin{align*}
-(c_0T + c_2P + c_3Q)T + (a_0T + a_2P + a_3Q)P + (b_0T + b_2P + b_3Q)Q &= kb_3(Q^2 - PR - TS),\\
-(c_1T + c_2Q + c_3R)T + (a_1T + a_2Q + a_3R)Q + (b_1T + b_2Q + b_3R)R &= -kb_1(Q^2 - PR - TS);
\end{align*}
and similarly, from the second set of equations, we obtain
\begin{align*}
-(d_0T + d_2P + d_3Q)T + (a_0T + a_2P + a_3Q)Q + (b_0T + b_2P + b_3Q)R &= -ka_3(Q^2 - PR - TS),\\
-(d_1T + d_2Q + d_3R)T + (a_1T + a_2Q + a_3R)Q + (b_1T + b_2Q + b_3R)R &= ka_2(Q^2 - PR - TS).
\end{align*}

\medskip

40.\ We have the inverse system
\[
(T, P, Q, R, S) = \frac{1}{k^2}
\begin{pmatrix}
(dc)_{01}, & (ad)_{31}, & 2(bd)_{31}, & (cb)_{01}, & (ab)_{01}\\
(dc)_{21}, & (ad)_{13}, & 2(bd)_{13}, & (cb)_{21}, & (ab)_{21}\\
(dc)_{31}, & (ad)_{31} - k, & 2(bd)_{31}, & (cb)_{31}, & (ab)_{31}\\
(dc)_{13}, & (ad)_{32}, & 2(bd)_{32}, & (cb)_{13}, & (ab)_{13}\\
(dc)_{32}, & (ad)_{02}, & 2(bd)_{02}, & (cb)_{32}, & (ab)_{32}
\end{pmatrix}
(T', P', Q', R', S'),
\]
read
\[
T' = (dc)_{01}T + (ad)_{31}P + 2(bd)_{31}Q + (cb)_{01}R + (ab)_{01}S;
\]
and in particular observe that, in the expression for $Q'$, the term containing $Q$ is $\{-k + 2(bd)_{13}\}Q$, where, in virtue of $(ac + bd)_{13} = k$, the coefficient of $Q$ is also $-(ac)_{13}$.

\medskip

41.\ These equations give
\begin{align*}
(a_0, a_2, a_3)(T', P', Q') &= k(d_3, ., -a_3, -b_3, b_0, b_2, b_3, ., ., c_3, a_3, b_3, ., d_3, c_3, .)(T, P, Q, R, S),\\
(a_1, a_2, a_3)(T', Q', R') &= k(-d_2, ., a_2, b_2, ., c_2, a_2, b_2, ., d_2, c_2, .)(T, P, Q, R, S).
\end{align*}

From the first set of equations, multiplying the first, third and fourth equations by $T$, $P$, $Q$ respectively, and adding, and again multiplying the second, third and fourth equations by $T$, $Q$, $R$ respectively, and adding, we deduce
\begin{align*}
&T(-a_0S - c_0R + d_0Q) + P(-a_2S - c_2R + d_2Q) + Q(-a_3S - c_3R + d_3Q) = -kc_3(Q^2 - PR - TS),\\
&T(-b_0S + c_0Q - d_0P) + P(-b_2S + c_2Q - d_2P) + Q(-b_3S + c_3Q - d_3P) = kd_3(Q^2 - PR - TS);
\end{align*}
and in like manner, from the second set of equations,
\begin{align*}
&T(-a_0S - c_0R + d_0Q) + Q(-a_2S - c_2R + d_2Q) + R(-a_3S - c_3R + d_3Q) = -kd_3(Q^2 - PR - TS),\\
&T(-b_0S + c_0Q - d_0P) + Q(-b_2S + c_2Q - d_2P) + R(-b_3S + c_3Q - d_3P) = kc_2(Q^2 - PR - TS).
\end{align*}

\medskip

42.\ Assume that $T$, $P$, $Q$, $R$, $S$ and $T'$, $P'$, $Q'$, $R'$, $S'$ are linearly connected as above; and write
\begin{align*}
T : P : Q : R : S &= 1 : A : H : B : H^2 - AB,\\
T' : P' : Q' : R' : S' &= 1 : A' : H' : B' : H'^2 - A'B',
\end{align*}
equations which establish a like relation between $1$, $A$, $H$, $B$, $H^2 - AB$, and $1$, $A'$, $H'$, $B'$, $H'^2 - A'B'$.

Observe that these forms are admissible since, if
\[
T^2 - PR - QS = 0,
\]
then also
\[
Q'^2 - P'R' - T'S' = 0.
\]

The quantities $A$, $H$, $B$ were taken as the parameters of a theta-function; viz.\ taking $A, H, B = A + i\alpha$, $H + i\eta$, $B + i\beta$, then $(\alpha, \eta, \beta \mid x, y)^2$ must be a positive form (or what is the same thing, $\alpha$ and $\alpha\beta - \eta^2$ must be positive). If $A'$, $H'$, $B'$ are also the parameters of a theta-function, then writing $A', H', B' = A' + i\alpha'$, $H' + i\eta'$, $B' + i\beta'$, $(\alpha', \eta', \beta' \mid x, y)^2$ must also be a positive form.

It can be shown that, $A'$, $H'$, $B'$ being determined as above, the former condition implies the latter one; viz.\ if the form $(\alpha, \eta, \beta \mid x, y)^2$ be positive, then also the form $(\alpha', \eta', \beta' \mid x, y)^2$ will be positive.

\medskip

43.\ Write
\begin{align*}
A &= A + i\alpha, & A' &= A' + i\alpha',\\
B &= B + i\beta, & B' &= B' + i\beta',\\
H &= H + i\eta, & H' &= H' + i\eta',
\end{align*}
and
\[
\Delta = H^2 - AB = (H'^2 - \alpha\beta) - (\eta^2 - \alpha'\beta'),
\]
and let the quintic matrix in the foregoing transformation
\[
(T', P', Q', R', S') = M(T, P, Q, R, S)
\]
be represented by
\[
\begin{pmatrix}
t, & p, & q, & r, & s\\
t_1, & p_1, & q_1, & r_1, & s_1\\
t_2, & p_2, & q_2, & r_2, & s_2\\
t_3, & p_3, & q_3, & r_3, & s_3\\
t_4, & p_4, & q_4, & r_4, & s_4
\end{pmatrix}.
\]

\newpage

44.\ It is to be shown that $\alpha'x^2 + 2\eta'xy + \beta'y^2$ is definite and positive.

We require $\alpha'$, $\eta'$, $\beta'$; we have
\begin{align*}
A' &= \frac{A't_1 + H't_2 + B't_3}{t}, & H' &= \frac{A'p_1 + H'p_2 + B'p_3}{p},\\
B' &= \frac{A'r_1 + H'r_2 + B'r_3}{r}, & \Delta' &= \frac{A's_1 + H's_2 + B's_3}{s},
\end{align*}
and thence
\[
\alpha' = \frac{\partial A'}{\partial i}, \quad \eta' = \frac{\partial H'}{\partial i}, \quad \beta' = \frac{\partial B'}{\partial i}.
\]

Now
\[
\alpha' = \frac{A't_1 + H't_2 + B't_3}{t} - \frac{(A't_1 + H't_2 + B't_3)^2}{t^2}\alpha,
\]
and so for $\eta'$, $\beta'$ with only the change of $t_1$, $t_2$, $t_3$ into $p_1$, $p_2$, $p_3$ and $r_1$, $r_2$, $r_3$ and $s_1$, $s_2$, $s_3$ respectively; we find, omitting a factor $\Delta$ throughout,
\[
\alpha' = \eta^2 - \alpha\beta, \quad \eta' = \eta\beta - \alpha\eta, \quad \beta' = \beta^2 - \eta^2.
\]

But we have identically
\[
(\eta^2 - \alpha\beta)x^2 + 2(\eta\beta - \alpha\eta)xy + (\beta^2 - \eta^2)y^2 = (\alpha x^2 + 2\eta xy + \beta y^2)(\eta x + \beta y)(\eta x - \alpha y),
\]
or, what is the same thing,
\[
\alpha' = \alpha\Delta, \quad \eta' = \eta\Delta, \quad \beta' = \beta\Delta.
\]

Hence $\alpha'x^2 + 2\eta'xy + \beta'y^2 = \Delta(\alpha x^2 + 2\eta xy + \beta y^2)$, which is positive if $\Delta$ is positive, that is, if $\alpha\beta - \eta^2$ is positive; and this is the required condition.

\bigskip

\subsection*{The Transformation. Art.~Nos.\ 45 to 60.}

\textbf{Art.~Nos.~45--48.]
}

45.\ I consider the transformation of the theta-functions with the parameters $(A, H, B)$ into theta-functions with the parameters $(A', H', B')$. Writing
\[
\vartheta_{\mu,\nu}(x, y \mid A, H, B) = \sum\sum (-1)^{\mu q + \nu r} \exp[i\pi\{(2m + \mu)x + (2n + \nu)y + (A, H, B \mid m, n)^2\}],
\]
and similarly
\[
\vartheta_{\mu',\nu'}(x', y' \mid A', H', B') = \sum\sum (-1)^{\mu' q + \nu' r} \exp[i\pi\{(2m' + \mu')x' + (2n' + \nu')y' + (A', H', B' \mid m', n')^2\}],
\]
the formula of transformation is
\[
\vartheta_{\mu',\nu'}(x', y') = C \cdot \vartheta_{\mu,\nu}(x, y),
\]
where $C$ is a constant depending on the characteristics and the matrix of transformation, and where the new arguments $x'$, $y'$ are linear functions of the original arguments $x$, $y$.

\medskip

46.\ The relation between the arguments is
\[
x' = \frac{(ab)_{01}x + (ab)_{21}y}{k}, \quad y' = \frac{(ac)_{01}x + (ac)_{21}y}{k},
\]
or, as we may write it,
\[
\begin{pmatrix} x' \\ y' \end{pmatrix}
= \frac{1}{k}
\begin{pmatrix} (ab)_{01}, & (ab)_{21} \\ (ac)_{01}, & (ac)_{21} \end{pmatrix}
\begin{pmatrix} x \\ y \end{pmatrix}.
\]

\medskip

47.\ The value of the constant $C$ may be written
\[
C = \frac{1}{\sqrt{k}} \cdot \frac{1}{\sqrt[4]{\Delta}} \cdot (-1)^{\frac{1}{2}(\mu q + \nu r + \mu' q + \nu' r)} \cdot \exp\left\{-\frac{i\pi}{k}\{(ab)_{01}x^2 + 2(ac)_{01}xy + (ad)_{01}y^2\}\right\},
\]
where $\Delta = (ac)_{01}(bd)_{01} - (ad)_{01}(bc)_{01}$.

\medskip

48.\ The foregoing formula may be verified by means of the equations which define the theta-functions; viz.\ substituting for $\vartheta_{\mu',\nu'}(x', y')$ its value in terms of $\vartheta_{\mu,\nu}(x, y)$, the equations
\begin{align*}
\vartheta_{\mu',\nu'}(x'+1, y') &= (-1)^{\mu' q}\vartheta_{\mu',\nu'}(x', y'),\\
\vartheta_{\mu',\nu'}(x', y'+1) &= (-1)^{\nu' r}\vartheta_{\mu',\nu'}(x', y'),\\
\vartheta_{\mu',\nu'}(x'+A', y'+H') &= (-1)^{q}\exp[-2i\pi(2x'+A')]\vartheta_{\mu',\nu'}(x', y'),\\
\vartheta_{\mu',\nu'}(x'+H', y'+B') &= (-1)^{r}\exp[-2i\pi(2y'+B')]\vartheta_{\mu',\nu'}(x', y'),
\end{align*}
should be respectively equivalent to the corresponding equations for $\vartheta_{\mu,\nu}(x, y)$.

\bigskip

\subsection*{Conclusion.}

The foregoing investigation completes the theory of the linear transformation of the double-theta functions. The results obtained are in substantial agreement with those of Hermite, but the notation has been modified and the demonstration rendered more systematic. The theory of the quintic matrix and its automorphic property in regard to the form $Q^2 - PR - TS$ presents features of considerable interest; and the arithmetical theory of the $k$-matrices, with the enumeration $1 + k + k^2 + k^3$ of the distinct forms, appears to be new.

\newpage

\end{document}
